% Options for packages loaded elsewhere
\PassOptionsToPackage{unicode}{hyperref}
\PassOptionsToPackage{hyphens}{url}
%
\documentclass[
  11pt,
  a4paper,
]{article}
\usepackage{amsmath,amssymb}
\usepackage{iftex}
\ifPDFTeX
  \usepackage[T1]{fontenc}
  \usepackage[utf8]{inputenc}
  \usepackage{textcomp} % provide euro and other symbols
\else % if luatex or xetex
  \usepackage{unicode-math} % this also loads fontspec
  \defaultfontfeatures{Scale=MatchLowercase}
  \defaultfontfeatures[\rmfamily]{Ligatures=TeX,Scale=1}
\fi
\usepackage{lmodern}
\ifPDFTeX\else
  % xetex/luatex font selection
  \setmainfont[]{Tinos}
  \setmonofont[]{DejaVu Sans Mono}
\fi
% Use upquote if available, for straight quotes in verbatim environments
\IfFileExists{upquote.sty}{\usepackage{upquote}}{}
\IfFileExists{microtype.sty}{% use microtype if available
  \usepackage[]{microtype}
  \UseMicrotypeSet[protrusion]{basicmath} % disable protrusion for tt fonts
}{}
\makeatletter
\@ifundefined{KOMAClassName}{% if non-KOMA class
  \IfFileExists{parskip.sty}{%
    \usepackage{parskip}
  }{% else
    \setlength{\parindent}{0pt}
    \setlength{\parskip}{6pt plus 2pt minus 1pt}}
}{% if KOMA class
  \KOMAoptions{parskip=half}}
\makeatother
\usepackage{xcolor}
\usepackage[top=2cm,bottom=2cm,left=2cm,right=2cm]{geometry}
\usepackage{color}
\usepackage{fancyvrb}
\newcommand{\VerbBar}{|}
\newcommand{\VERB}{\Verb[commandchars=\\\{\}]}
\DefineVerbatimEnvironment{Highlighting}{Verbatim}{commandchars=\\\{\}}
% Add ',fontsize=\small' for more characters per line
\newenvironment{Shaded}{}{}
\newcommand{\AlertTok}[1]{\textcolor[rgb]{1.00,0.00,0.00}{\textbf{#1}}}
\newcommand{\AnnotationTok}[1]{\textcolor[rgb]{0.38,0.63,0.69}{\textbf{\textit{#1}}}}
\newcommand{\AttributeTok}[1]{\textcolor[rgb]{0.49,0.56,0.16}{#1}}
\newcommand{\BaseNTok}[1]{\textcolor[rgb]{0.25,0.63,0.44}{#1}}
\newcommand{\BuiltInTok}[1]{\textcolor[rgb]{0.00,0.50,0.00}{#1}}
\newcommand{\CharTok}[1]{\textcolor[rgb]{0.25,0.44,0.63}{#1}}
\newcommand{\CommentTok}[1]{\textcolor[rgb]{0.38,0.63,0.69}{\textit{#1}}}
\newcommand{\CommentVarTok}[1]{\textcolor[rgb]{0.38,0.63,0.69}{\textbf{\textit{#1}}}}
\newcommand{\ConstantTok}[1]{\textcolor[rgb]{0.53,0.00,0.00}{#1}}
\newcommand{\ControlFlowTok}[1]{\textcolor[rgb]{0.00,0.44,0.13}{\textbf{#1}}}
\newcommand{\DataTypeTok}[1]{\textcolor[rgb]{0.56,0.13,0.00}{#1}}
\newcommand{\DecValTok}[1]{\textcolor[rgb]{0.25,0.63,0.44}{#1}}
\newcommand{\DocumentationTok}[1]{\textcolor[rgb]{0.73,0.13,0.13}{\textit{#1}}}
\newcommand{\ErrorTok}[1]{\textcolor[rgb]{1.00,0.00,0.00}{\textbf{#1}}}
\newcommand{\ExtensionTok}[1]{#1}
\newcommand{\FloatTok}[1]{\textcolor[rgb]{0.25,0.63,0.44}{#1}}
\newcommand{\FunctionTok}[1]{\textcolor[rgb]{0.02,0.16,0.49}{#1}}
\newcommand{\ImportTok}[1]{\textcolor[rgb]{0.00,0.50,0.00}{\textbf{#1}}}
\newcommand{\InformationTok}[1]{\textcolor[rgb]{0.38,0.63,0.69}{\textbf{\textit{#1}}}}
\newcommand{\KeywordTok}[1]{\textcolor[rgb]{0.00,0.44,0.13}{\textbf{#1}}}
\newcommand{\NormalTok}[1]{#1}
\newcommand{\OperatorTok}[1]{\textcolor[rgb]{0.40,0.40,0.40}{#1}}
\newcommand{\OtherTok}[1]{\textcolor[rgb]{0.00,0.44,0.13}{#1}}
\newcommand{\PreprocessorTok}[1]{\textcolor[rgb]{0.74,0.48,0.00}{#1}}
\newcommand{\RegionMarkerTok}[1]{#1}
\newcommand{\SpecialCharTok}[1]{\textcolor[rgb]{0.25,0.44,0.63}{#1}}
\newcommand{\SpecialStringTok}[1]{\textcolor[rgb]{0.73,0.40,0.53}{#1}}
\newcommand{\StringTok}[1]{\textcolor[rgb]{0.25,0.44,0.63}{#1}}
\newcommand{\VariableTok}[1]{\textcolor[rgb]{0.10,0.09,0.49}{#1}}
\newcommand{\VerbatimStringTok}[1]{\textcolor[rgb]{0.25,0.44,0.63}{#1}}
\newcommand{\WarningTok}[1]{\textcolor[rgb]{0.38,0.63,0.69}{\textbf{\textit{#1}}}}
\usepackage{longtable,booktabs,array}
\usepackage{calc} % for calculating minipage widths
% Correct order of tables after \paragraph or \subparagraph
\usepackage{etoolbox}
\makeatletter
\patchcmd\longtable{\par}{\if@noskipsec\mbox{}\fi\par}{}{}
\makeatother
% Allow footnotes in longtable head/foot
\IfFileExists{footnotehyper.sty}{\usepackage{footnotehyper}}{\usepackage{footnote}}
\makesavenoteenv{longtable}
\setlength{\emergencystretch}{3em} % prevent overfull lines
\providecommand{\tightlist}{%
  \setlength{\itemsep}{0pt}\setlength{\parskip}{0pt}}
\setcounter{secnumdepth}{-\maxdimen} % remove section numbering
\usepackage{fancyhdr}
\usepackage{etoolbox}
\AtBeginEnvironment{longtable}{\setlength{\parskip}{0pt}}
\pagestyle{fancy}
\fancyhf{}
\setlength{\headheight}{16pt}
\lhead{\footnotesize\textbf{DSAI1003 -- Fundamental Programming Concepts in Python}}
\rhead{\footnotesize\textbf{Bài 3: Câu lệnh rẽ nhánh}}
\lfoot{\scriptsize Khoa Khoa học Dữ liệu và Trí tuệ Nhân tạo -- Đại học Kinh tế Quốc dân}
\rfoot{\footnotesize Trang \thepage}
\renewcommand{\headrulewidth}{0.4pt}
\renewcommand{\footrulewidth}{0.3pt}
\setlength{\emergencystretch}{3em}
\ifLuaTeX
  \usepackage{selnolig}  % disable illegal ligatures
\fi
\usepackage{bookmark}
\IfFileExists{xurl.sty}{\usepackage{xurl}}{} % add URL line breaks if available
\urlstyle{same}
\hypersetup{
  hidelinks,
  pdfcreator={LaTeX via pandoc}}

\title{Bài 3: Câu lệnh rẽ nhánh}
\author{Khoa Khoa học Dữ liệu và Trí tuệ Nhân tạo (FDA) - Đại học Kinh tế Quốc dân}
\date{Mã học phần: DSAI1003}

\begin{document}
\maketitle

{
\setcounter{tocdepth}{3}
\tableofcontents
}
\section{Bài 3. Câu lệnh rẽ
nhánh}\label{buxe0i-3.-cuxe2u-lux1ec7nh-rux1ebd-nhuxe1nh}

\textbf{Cập nhật lần cuối:} 19/09/2026

\begin{quote}
\textbf{Nguồn:} Chương 3 -- \emph{Decisions}, trong giáo trình
\emph{Python for Everyone} của Cay Horstmann và Rance Necaise (2/e).

Bài giảng này giữ cấu trúc, thuật ngữ, ví dụ và trọng tâm sư phạm của
Chương 3, đồng thời sắp xếp lại nội dung để thuận tiện cho giảng dạy
trực tiếp trên lớp. Các ví dụ vẫn theo phong cách đặt tên biến và định
dạng chuỗi bằng toán tử \texttt{\%} như trong giáo trình.
\end{quote}

\begin{center}\rule{0.5\linewidth}{0.5pt}\end{center}

\subsection{Giới thiệu bài
học}\label{giux1edbi-thiux1ec7u-buxe0i-hux1ecdc}

Ở Bài 2, phần lớn các chương trình được thực hiện theo thứ tự từ trên
xuống dưới. Cách thực thi này phù hợp với các phép tính trực tiếp, nhưng
các chương trình thực tế còn cần khả năng \textbf{đưa ra quyết định}.

Một quyết định cho phép chương trình chọn hành động khác nhau tùy theo
dữ liệu hoặc hoàn cảnh. Ví dụ, chương trình có thể cần:

\begin{itemize}
\tightlist
\item
  áp dụng mức giảm giá khác nhau cho đơn hàng nhỏ và đơn hàng lớn;
\item
  từ chối số tầng không hợp lệ trong chương trình mô phỏng thang máy;
\item
  phân loại động đất theo độ lớn;
\item
  kiểm tra một chuỗi có đúng định dạng yêu cầu hay không;
\item
  chọn công thức tính thuế theo tình trạng hôn nhân và thu nhập;
\item
  kiểm tra đồng thời nhiều điều kiện.
\end{itemize}

Ý tưởng trung tâm của bài học là:

\begin{Shaded}
\begin{Highlighting}[]
\NormalTok{Đánh giá một điều kiện}
\NormalTok{          ↓}
\NormalTok{    Đúng hay Sai?}
\NormalTok{       ↙      ↘}
\NormalTok{  Hành động A  Hành động B}
\NormalTok{       ↘      ↙}
\NormalTok{  Tiếp tục chương trình}
\end{Highlighting}
\end{Shaded}

Nội dung được phát triển từ câu lệnh \texttt{if} hai nhánh, sau đó mở
rộng sang rẽ nhánh lồng nhau, nhiều phương án, biểu thức Boolean, phân
tích chuỗi, thiết kế ca kiểm thử và kiểm tra dữ liệu đầu vào.

\begin{center}\rule{0.5\linewidth}{0.5pt}\end{center}

\subsection{Kiến thức và kỹ năng cần
đạt}\label{kiux1ebfn-thux1ee9c-vuxe0-kux1ef9-nux103ng-cux1ea7n-ux111ux1ea1t}

Sau khi hoàn thành bài học này, sinh viên có thể:

\begin{itemize}
\tightlist
\item
  Sử dụng các cấu trúc \texttt{if}, \texttt{if/else} và
  \texttt{if/elif/else}.
\item
  Giải thích vai trò của \textbf{điều kiện} trong luồng điều khiển của
  chương trình.
\item
  Sử dụng thụt lề đúng để xác định các khối lệnh.
\item
  So sánh số và chuỗi bằng các toán tử quan hệ.
\item
  Phân biệt phép gán \texttt{=} với phép kiểm tra bằng nhau \texttt{==}.
\item
  Tránh so sánh bằng tuyệt đối với các số thực được tạo ra từ phép tính
  khi cần kiểm tra gần bằng.
\item
  Thiết kế quyết định hai nhánh từ mô tả bài toán.
\item
  Cài đặt \textbf{rẽ nhánh lồng nhau} cho các quyết định nhiều tầng.
\item
  Sử dụng \texttt{elif} cho các phương án loại trừ lẫn nhau.
\item
  Sắp xếp điều kiện theo thứ tự phù hợp, đặc biệt với các ngưỡng chồng
  lấn.
\item
  Đọc và xây dựng \textbf{lưu đồ} đơn giản.
\item
  Thiết kế \textbf{ca kiểm thử} bao phủ các nhánh và giá trị biên.
\item
  Sử dụng hai giá trị Boolean \texttt{True} và \texttt{False}.
\item
  Kết hợp điều kiện bằng \texttt{and}, \texttt{or} và \texttt{not}.
\item
  Giải thích thứ tự ưu tiên và cơ chế đánh giá ngắn mạch.
\item
  Áp dụng định luật De Morgan để biến đổi biểu thức logic.
\item
  Phân tích chuỗi bằng toán tử thành viên và các phương thức chuỗi.
\item
  Kiểm tra tính hợp lệ của dữ liệu đầu vào trước khi xử lý.
\item
  Chuẩn hóa dữ liệu văn bản bằng \texttt{upper()} hoặc \texttt{lower()}
  khi không cần phân biệt chữ hoa/chữ thường.
\end{itemize}

\begin{center}\rule{0.5\linewidth}{0.5pt}\end{center}

\subsection{Cấu trúc bài học}\label{cux1ea5u-truxfac-buxe0i-hux1ecdc}

\begin{enumerate}
\def\labelenumi{\arabic{enumi}.}
\tightlist
\item
  Câu lệnh \texttt{if}
\item
  Các toán tử quan hệ
\item
  Rẽ nhánh lồng nhau
\item
  Nhiều phương án
\item
  Giải quyết vấn đề: Lưu đồ
\item
  Giải quyết vấn đề: Ca kiểm thử
\item
  Biến Boolean và toán tử Boolean
\item
  Phân tích chuỗi
\item
  Ứng dụng: Kiểm tra dữ liệu đầu vào
\item
  Các chủ đề và Toolbox tùy chọn
\item
  Tổng kết và bài tập
\end{enumerate}

\begin{center}\rule{0.5\linewidth}{0.5pt}\end{center}

\section{\texorpdfstring{3.1 Câu lệnh
\texttt{if}}{3.1 Câu lệnh if}}\label{cuxe2u-lux1ec7nh-if}

Chương trình thường cần thực hiện một hành động khi điều kiện đúng và
một hành động khác khi điều kiện sai.

Câu lệnh \texttt{if} cho phép chương trình thực hiện việc rẽ nhánh này.

\subsection{3.1.1 Quyết định hai
nhánh}\label{quyux1ebft-ux111ux1ecbnh-hai-nhuxe1nh}

Xét hệ thống thang máy trong một tòa nhà bỏ qua số tầng 13. Khi người
dùng chọn tầng hiển thị lớn hơn 13, tầng vật lý thực tế thấp hơn 1 đơn
vị. Ví dụ, tầng hiển thị 20 tương ứng với tầng vật lý 19.

Ta có thể viết:

\begin{Shaded}
\begin{Highlighting}[]
\ControlFlowTok{if}\NormalTok{ floor }\OperatorTok{\textgreater{}} \DecValTok{13}\NormalTok{ :}
\NormalTok{    actualFloor }\OperatorTok{=}\NormalTok{ floor }\OperatorTok{{-}} \DecValTok{1}
\ControlFlowTok{else}\NormalTok{ :}
\NormalTok{    actualFloor }\OperatorTok{=}\NormalTok{ floor}
\end{Highlighting}
\end{Shaded}

Về mặt ý tưởng:

\begin{Shaded}
\begin{Highlighting}[]
\NormalTok{                 floor \textgreater{} 13 ?}
\NormalTok{                  /       \textbackslash{}}
\NormalTok{              Đúng        Sai}
\NormalTok{               /            \textbackslash{}}
\NormalTok{ actualFloor = floor {-} 1   actualFloor = floor}
\end{Highlighting}
\end{Shaded}

Chỉ \textbf{một nhánh} được thực hiện.

Nếu \texttt{floor\ \textgreater{}\ 13} là đúng, Python thực hiện khối
lệnh đầu tiên và bỏ qua khối \texttt{else}. Nếu điều kiện sai, Python bỏ
qua khối đầu và thực hiện nhánh \texttt{else}.

\begin{center}\rule{0.5\linewidth}{0.5pt}\end{center}

\subsection{3.1.2 Cú pháp tổng
quát}\label{cuxfa-phuxe1p-tux1ed5ng-quuxe1t}

\begin{Shaded}
\begin{Highlighting}[]
\ControlFlowTok{if}\NormalTok{ condition :}
\NormalTok{    statements1}
\ControlFlowTok{else}\NormalTok{ :}
\NormalTok{    statements2}
\end{Highlighting}
\end{Shaded}

Nhánh \texttt{else} là tùy chọn:

\begin{Shaded}
\begin{Highlighting}[]
\ControlFlowTok{if}\NormalTok{ condition :}
\NormalTok{    statements}
\end{Highlighting}
\end{Shaded}

Ví dụ:

\begin{Shaded}
\begin{Highlighting}[]
\NormalTok{actualFloor }\OperatorTok{=}\NormalTok{ floor}

\ControlFlowTok{if}\NormalTok{ floor }\OperatorTok{\textgreater{}} \DecValTok{13}\NormalTok{ :}
\NormalTok{    actualFloor }\OperatorTok{=}\NormalTok{ actualFloor }\OperatorTok{{-}} \DecValTok{1}
\end{Highlighting}
\end{Shaded}

Cách viết này phù hợp khi trường hợp sai không cần hành động đặc biệt
nào.

\begin{center}\rule{0.5\linewidth}{0.5pt}\end{center}

\subsection{3.1.3 Câu lệnh ghép và khối
lệnh}\label{cuxe2u-lux1ec7nh-ghuxe9p-vuxe0-khux1ed1i-lux1ec7nh}

Câu lệnh \texttt{if} là một \textbf{câu lệnh ghép}. Dòng đầu của nó kết
thúc bằng dấu hai chấm \texttt{:}, và các câu lệnh thuộc cùng nhánh phải
được thụt lề.

\begin{Shaded}
\begin{Highlighting}[]
\ControlFlowTok{if}\NormalTok{ totalSales }\OperatorTok{\textgreater{}} \FloatTok{100.0}\NormalTok{ :}
\NormalTok{    discount }\OperatorTok{=}\NormalTok{ totalSales }\OperatorTok{*} \FloatTok{0.05}
\NormalTok{    totalSales }\OperatorTok{=}\NormalTok{ totalSales }\OperatorTok{{-}}\NormalTok{ discount}
    \BuiltInTok{print}\NormalTok{(}\StringTok{"You received a discount of"}\NormalTok{, discount)}
\end{Highlighting}
\end{Shaded}

Ba dòng được thụt lề tạo thành một khối lệnh.

Các quy tắc quan trọng:

\begin{itemize}
\tightlist
\item
  Dòng tiêu đề kết thúc bằng \texttt{:}.
\item
  Các câu lệnh trong cùng một khối phải thụt lề ở cùng mức.
\item
  \texttt{if} và \texttt{else} tương ứng phải thẳng hàng.
\item
  Khối lệnh kết thúc khi mức thụt lề quay lại mức bên ngoài.
\end{itemize}

Trong Python, thụt lề là một phần của cú pháp, không chỉ để trình bày
đẹp hơn.

\begin{center}\rule{0.5\linewidth}{0.5pt}\end{center}

\subsection{Ví dụ -- Mô phỏng thang
máy}\label{vuxed-dux1ee5-muxf4-phux1ecfng-thang-muxe1y}

\begin{Shaded}
\begin{Highlighting}[]
\CommentTok{\#\#}
\CommentTok{\# This program simulates an elevator panel that skips floor 13.}
\CommentTok{\#}

\NormalTok{floor }\OperatorTok{=} \BuiltInTok{int}\NormalTok{(}\BuiltInTok{input}\NormalTok{(}\StringTok{"Floor: "}\NormalTok{))}

\ControlFlowTok{if}\NormalTok{ floor }\OperatorTok{\textgreater{}} \DecValTok{13}\NormalTok{ :}
\NormalTok{    actualFloor }\OperatorTok{=}\NormalTok{ floor }\OperatorTok{{-}} \DecValTok{1}
\ControlFlowTok{else}\NormalTok{ :}
\NormalTok{    actualFloor }\OperatorTok{=}\NormalTok{ floor}

\BuiltInTok{print}\NormalTok{(}\StringTok{"The elevator will travel to the actual floor"}\NormalTok{, actualFloor)}
\end{Highlighting}
\end{Shaded}

Ví dụ chạy chương trình:

\begin{Shaded}
\begin{Highlighting}[]
\NormalTok{Floor: 20}
\NormalTok{The elevator will travel to the actual floor 19}
\end{Highlighting}
\end{Shaded}

\begin{center}\rule{0.5\linewidth}{0.5pt}\end{center}

\section{Lỗi thường gặp 3.1 -- Tab và thụt lề không nhất
quán}\label{lux1ed7i-thux1b0ux1eddng-gux1eb7p-3.1-tab-vuxe0-thux1ee5t-lux1ec1-khuxf4ng-nhux1ea5t-quuxe1n}

Python yêu cầu các câu lệnh trong cùng một khối phải có mức thụt lề nhất
quán.

Một đoạn mã nhìn có vẻ thẳng hàng trong trình soạn thảo vẫn có thể gặp
lỗi nếu trộn ký tự tab và khoảng trắng.

Nên cấu hình trình soạn thảo để phím Tab chèn khoảng trắng. Một quy ước
phổ biến là dùng bốn khoảng trắng cho mỗi mức thụt lề:

\begin{Shaded}
\begin{Highlighting}[]
\ControlFlowTok{if}\NormalTok{ score }\OperatorTok{\textgreater{}=} \DecValTok{50}\NormalTok{ :}
    \BuiltInTok{print}\NormalTok{(}\StringTok{"Pass"}\NormalTok{)}
    \BuiltInTok{print}\NormalTok{(}\StringTok{"Continue to the next course"}\NormalTok{)}
\end{Highlighting}
\end{Shaded}

Điều quan trọng nhất là \textbf{tính nhất quán}.

\begin{center}\rule{0.5\linewidth}{0.5pt}\end{center}

\section{Mẹo lập trình 3.1 -- Tránh lặp mã trong các
nhánh}\label{mux1eb9o-lux1eadp-truxecnh-3.1-truxe1nh-lux1eb7p-muxe3-trong-cuxe1c-nhuxe1nh}

Xét đoạn mã:

\begin{Shaded}
\begin{Highlighting}[]
\ControlFlowTok{if}\NormalTok{ floor }\OperatorTok{\textgreater{}} \DecValTok{13}\NormalTok{ :}
\NormalTok{    actualFloor }\OperatorTok{=}\NormalTok{ floor }\OperatorTok{{-}} \DecValTok{1}
    \BuiltInTok{print}\NormalTok{(}\StringTok{"Actual floor:"}\NormalTok{, actualFloor)}
\ControlFlowTok{else}\NormalTok{ :}
\NormalTok{    actualFloor }\OperatorTok{=}\NormalTok{ floor}
    \BuiltInTok{print}\NormalTok{(}\StringTok{"Actual floor:"}\NormalTok{, actualFloor)}
\end{Highlighting}
\end{Shaded}

Câu lệnh \texttt{print} bị lặp lại ở cả hai nhánh.

Cách viết tốt hơn:

\begin{Shaded}
\begin{Highlighting}[]
\ControlFlowTok{if}\NormalTok{ floor }\OperatorTok{\textgreater{}} \DecValTok{13}\NormalTok{ :}
\NormalTok{    actualFloor }\OperatorTok{=}\NormalTok{ floor }\OperatorTok{{-}} \DecValTok{1}
\ControlFlowTok{else}\NormalTok{ :}
\NormalTok{    actualFloor }\OperatorTok{=}\NormalTok{ floor}

\BuiltInTok{print}\NormalTok{(}\StringTok{"Actual floor:"}\NormalTok{, actualFloor)}
\end{Highlighting}
\end{Shaded}

Đưa phần xử lý chung ra khỏi các nhánh giúp:

\begin{itemize}
\tightlist
\item
  giảm lặp mã;
\item
  chương trình ngắn hơn;
\item
  giảm nguy cơ sửa một bản nhưng quên sửa bản còn lại.
\end{itemize}

\begin{center}\rule{0.5\linewidth}{0.5pt}\end{center}

\section{Chủ đề bổ sung 3.1 -- Biểu thức điều
kiện}\label{chux1ee7-ux111ux1ec1-bux1ed5-sung-3.1-biux1ec3u-thux1ee9c-ux111iux1ec1u-kiux1ec7n}

Python có thể biểu diễn một lựa chọn hai trường hợp dưới dạng biểu thức:

\begin{Shaded}
\begin{Highlighting}[]
\NormalTok{actualFloor }\OperatorTok{=}\NormalTok{ floor }\OperatorTok{{-}} \DecValTok{1} \ControlFlowTok{if}\NormalTok{ floor }\OperatorTok{\textgreater{}} \DecValTok{13} \ControlFlowTok{else}\NormalTok{ floor}
\end{Highlighting}
\end{Shaded}

Tương đương với:

\begin{Shaded}
\begin{Highlighting}[]
\ControlFlowTok{if}\NormalTok{ floor }\OperatorTok{\textgreater{}} \DecValTok{13}\NormalTok{ :}
\NormalTok{    actualFloor }\OperatorTok{=}\NormalTok{ floor }\OperatorTok{{-}} \DecValTok{1}
\ControlFlowTok{else}\NormalTok{ :}
\NormalTok{    actualFloor }\OperatorTok{=}\NormalTok{ floor}
\end{Highlighting}
\end{Shaded}

Ở giai đoạn đầu học lập trình, dạng \texttt{if/else} đầy đủ thường dễ
đọc hơn, đặc biệt khi mỗi nhánh có nhiều hơn một hành động.

\begin{center}\rule{0.5\linewidth}{0.5pt}\end{center}

\subsection{Tự kiểm tra 3.1}\label{tux1ef1-kiux1ec3m-tra-3.1}

\subsubsection{Câu 1}\label{cuxe2u-1}

Giá trị của \texttt{actualFloor} là bao nhiêu khi \texttt{floor\ =\ 18}?

\begin{Shaded}
\begin{Highlighting}[]
\ControlFlowTok{if}\NormalTok{ floor }\OperatorTok{\textgreater{}} \DecValTok{13}\NormalTok{ :}
\NormalTok{    actualFloor }\OperatorTok{=}\NormalTok{ floor }\OperatorTok{{-}} \DecValTok{1}
\ControlFlowTok{else}\NormalTok{ :}
\NormalTok{    actualFloor }\OperatorTok{=}\NormalTok{ floor}
\end{Highlighting}
\end{Shaded}

\textbf{Đáp án}

\texttt{17}.

\subsubsection{Câu 2}\label{cuxe2u-2}

Giá trị của \texttt{actualFloor} là bao nhiêu khi \texttt{floor\ =\ 9}?

\textbf{Đáp án}

\texttt{9}.

\subsubsection{Câu 3}\label{cuxe2u-3}

Tại sao phải có dấu \texttt{:} sau điều kiện của \texttt{if}?

\textbf{Đáp án}

Dấu \texttt{:} đánh dấu dòng tiêu đề của câu lệnh ghép. Khối lệnh được
thụt lề ngay sau đó thuộc về dòng tiêu đề này.

\subsubsection{Câu 4}\label{cuxe2u-4}

Có thể cải thiện đoạn mã sau như thế nào?

\begin{Shaded}
\begin{Highlighting}[]
\ControlFlowTok{if}\NormalTok{ balance }\OperatorTok{\textless{}} \DecValTok{0}\NormalTok{ :}
\NormalTok{    status }\OperatorTok{=} \StringTok{"Overdrawn"}
    \BuiltInTok{print}\NormalTok{(status)}
\ControlFlowTok{else}\NormalTok{ :}
\NormalTok{    status }\OperatorTok{=} \StringTok{"OK"}
    \BuiltInTok{print}\NormalTok{(status)}
\end{Highlighting}
\end{Shaded}

\textbf{Đáp án}

Đưa câu lệnh \texttt{print} ra ngoài hai nhánh:

\begin{Shaded}
\begin{Highlighting}[]
\ControlFlowTok{if}\NormalTok{ balance }\OperatorTok{\textless{}} \DecValTok{0}\NormalTok{ :}
\NormalTok{    status }\OperatorTok{=} \StringTok{"Overdrawn"}
\ControlFlowTok{else}\NormalTok{ :}
\NormalTok{    status }\OperatorTok{=} \StringTok{"OK"}

\BuiltInTok{print}\NormalTok{(status)}
\end{Highlighting}
\end{Shaded}

\begin{center}\rule{0.5\linewidth}{0.5pt}\end{center}

\section{3.2 Các toán tử quan
hệ}\label{cuxe1c-touxe1n-tux1eed-quan-hux1ec7}

Một câu lệnh \texttt{if} cần một điều kiện có giá trị là \texttt{True}
hoặc \texttt{False}.

Nhiều điều kiện được tạo bằng cách so sánh hai giá trị với \textbf{toán
tử quan hệ}.

\subsection{3.2.1 Các toán tử quan
hệ}\label{cuxe1c-touxe1n-tux1eed-quan-hux1ec7-1}

\begin{longtable}[]{@{}ll@{}}
\toprule\noalign{}
Toán tử Python & Ý nghĩa \\
\midrule\noalign{}
\endhead
\bottomrule\noalign{}
\endlastfoot
\texttt{\textgreater{}} & Lớn hơn \\
\texttt{\textgreater{}=} & Lớn hơn hoặc bằng \\
\texttt{\textless{}} & Nhỏ hơn \\
\texttt{\textless{}=} & Nhỏ hơn hoặc bằng \\
\texttt{==} & Bằng \\
\texttt{!=} & Khác \\
\end{longtable}

Ví dụ:

\begin{Shaded}
\begin{Highlighting}[]
\DecValTok{3} \OperatorTok{\textless{}} \DecValTok{4}
\DecValTok{4} \OperatorTok{\textless{}=} \DecValTok{4}
\DecValTok{7} \OperatorTok{!=} \DecValTok{5}
\DecValTok{10} \OperatorTok{==} \DecValTok{2} \OperatorTok{*} \DecValTok{5}
\end{Highlighting}
\end{Shaded}

Mỗi biểu thức trên cho kết quả là một giá trị Boolean.

\begin{center}\rule{0.5\linewidth}{0.5pt}\end{center}

\subsection{\texorpdfstring{3.2.2 Phép gán \texttt{=} và phép kiểm tra
bằng
\texttt{==}}{3.2.2 Phép gán = và phép kiểm tra bằng ==}}\label{phuxe9p-guxe1n-vuxe0-phuxe9p-kiux1ec3m-tra-bux1eb1ng}

Hai toán tử này có chức năng hoàn toàn khác nhau.

Phép gán:

\begin{Shaded}
\begin{Highlighting}[]
\NormalTok{floor }\OperatorTok{=} \DecValTok{13}
\end{Highlighting}
\end{Shaded}

Phép kiểm tra bằng nhau:

\begin{Shaded}
\begin{Highlighting}[]
\ControlFlowTok{if}\NormalTok{ floor }\OperatorTok{==} \DecValTok{13}\NormalTok{ :}
    \BuiltInTok{print}\NormalTok{(}\StringTok{"Floor 13 selected"}\NormalTok{)}
\end{Highlighting}
\end{Shaded}

Có thể đọc như sau:

\begin{Shaded}
\begin{Highlighting}[]
\NormalTok{=   → gán}
\NormalTok{==  → có bằng nhau không?}
\end{Highlighting}
\end{Shaded}

\begin{center}\rule{0.5\linewidth}{0.5pt}\end{center}

\subsection{3.2.3 So sánh chuỗi}\label{so-suxe1nh-chuux1ed7i}

Chuỗi cũng có thể được so sánh.

\begin{Shaded}
\begin{Highlighting}[]
\NormalTok{name1 }\OperatorTok{=} \StringTok{"John Wayne"}
\NormalTok{name2 }\OperatorTok{=} \StringTok{"John Wayne"}

\ControlFlowTok{if}\NormalTok{ name1 }\OperatorTok{==}\NormalTok{ name2 :}
    \BuiltInTok{print}\NormalTok{(}\StringTok{"The strings are identical."}\NormalTok{)}
\end{Highlighting}
\end{Shaded}

So sánh chuỗi phân biệt chữ hoa và chữ thường:

\begin{Shaded}
\begin{Highlighting}[]
\CommentTok{"John"} \OperatorTok{==} \StringTok{"john"}
\end{Highlighting}
\end{Shaded}

cho kết quả \texttt{False}.

Để hai chuỗi bằng nhau, mọi ký tự và vị trí của ký tự phải trùng khớp.

\begin{center}\rule{0.5\linewidth}{0.5pt}\end{center}

\subsection{3.2.4 Thứ tự ưu tiên toán
tử}\label{thux1ee9-tux1ef1-ux1b0u-tiuxean-touxe1n-tux1eed}

Các toán tử số học có độ ưu tiên cao hơn toán tử quan hệ.

\begin{Shaded}
\begin{Highlighting}[]
\NormalTok{floor }\OperatorTok{{-}} \DecValTok{1} \OperatorTok{\textless{}} \DecValTok{13}
\end{Highlighting}
\end{Shaded}

Python tính trước:

\begin{Shaded}
\begin{Highlighting}[]
\NormalTok{floor {-} 1}
\end{Highlighting}
\end{Shaded}

sau đó mới so sánh kết quả với \texttt{13}.

Với các biểu thức phức tạp, nên dùng dấu ngoặc để tăng tính dễ đọc.

\begin{center}\rule{0.5\linewidth}{0.5pt}\end{center}

\section{Lỗi thường gặp 3.2 -- So sánh chính xác số thực dấu phẩy
động}\label{lux1ed7i-thux1b0ux1eddng-gux1eb7p-3.2-so-suxe1nh-chuxednh-xuxe1c-sux1ed1-thux1ef1c-dux1ea5u-phux1ea9y-ux111ux1ed9ng}

Các phép tính với \texttt{float} có thể xuất hiện sai số làm tròn rất
nhỏ.

Ví dụ:

\begin{Shaded}
\begin{Highlighting}[]
\ImportTok{from}\NormalTok{ math }\ImportTok{import}\NormalTok{ sqrt}

\NormalTok{x }\OperatorTok{=}\NormalTok{ sqrt(}\FloatTok{2.0}\NormalTok{)}
\BuiltInTok{print}\NormalTok{(x }\OperatorTok{*}\NormalTok{ x)}
\end{Highlighting}
\end{Shaded}

Kết quả có thể rất gần \texttt{2.0}, nhưng biểu diễn bên trong không
nhất thiết chính xác tuyệt đối bằng \texttt{2.0}.

Vì vậy phép kiểm tra sau có thể không phù hợp:

\begin{Shaded}
\begin{Highlighting}[]
\ControlFlowTok{if}\NormalTok{ x }\OperatorTok{*}\NormalTok{ x }\OperatorTok{==} \FloatTok{2.0}\NormalTok{ :}
    \BuiltInTok{print}\NormalTok{(}\StringTok{"Exactly two"}\NormalTok{)}
\end{Highlighting}
\end{Shaded}

Khi cần kiểm tra gần bằng, so sánh độ chênh với một ngưỡng nhỏ:

\begin{Shaded}
\begin{Highlighting}[]
\NormalTok{EPSILON }\OperatorTok{=} \FloatTok{1E{-}14}

\ControlFlowTok{if} \BuiltInTok{abs}\NormalTok{(x }\OperatorTok{*}\NormalTok{ x }\OperatorTok{{-}} \FloatTok{2.0}\NormalTok{) }\OperatorTok{\textless{}}\NormalTok{ EPSILON :}
    \BuiltInTok{print}\NormalTok{(}\StringTok{"Approximately two"}\NormalTok{)}
\end{Highlighting}
\end{Shaded}

Mẫu tổng quát:

\begin{Shaded}
\begin{Highlighting}[]
\ControlFlowTok{if} \BuiltInTok{abs}\NormalTok{(x }\OperatorTok{{-}}\NormalTok{ y) }\OperatorTok{\textless{}}\NormalTok{ epsilon :}
    \CommentTok{\# Coi x và y là gần bằng nhau.}
\end{Highlighting}
\end{Shaded}

Giá trị \texttt{epsilon} phải được chọn phù hợp với thang đo và mục đích
của phép tính.

\begin{center}\rule{0.5\linewidth}{0.5pt}\end{center}

\section{Chủ đề bổ sung 3.2 -- Thứ tự từ điển của
chuỗi}\label{chux1ee7-ux111ux1ec1-bux1ed5-sung-3.2-thux1ee9-tux1ef1-tux1eeb-ux111iux1ec3n-cux1ee7a-chuux1ed7i}

Các toán tử như \texttt{\textless{}} và \texttt{\textgreater{}} cũng có
thể dùng với chuỗi.

Python so sánh chuỗi theo thứ tự từ điển dựa trên mã ký tự.

Ví dụ, chữ hoa/chữ thường có ảnh hưởng:

\begin{Shaded}
\begin{Highlighting}[]
\CommentTok{"John"} \OperatorTok{\textless{}} \StringTok{"john"}
\end{Highlighting}
\end{Shaded}

Kết quả dựa trên thứ tự mã ký tự, không phải quy tắc từ điển bỏ qua chữ
hoa/chữ thường.

Nếu muốn so sánh không phân biệt hoa thường, có thể chuẩn hóa trước:

\begin{Shaded}
\begin{Highlighting}[]
\NormalTok{name1 }\OperatorTok{=}\NormalTok{ name1.lower()}
\NormalTok{name2 }\OperatorTok{=}\NormalTok{ name2.lower()}
\end{Highlighting}
\end{Shaded}

sau đó mới so sánh.

\begin{center}\rule{0.5\linewidth}{0.5pt}\end{center}

\section{\texorpdfstring{CÁCH LÀM 3.1 -- Cài đặt một câu lệnh
\texttt{if}}{CÁCH LÀM 3.1 -- Cài đặt một câu lệnh if}}\label{cuxe1ch-luxe0m-3.1-cuxe0i-ux111ux1eb7t-mux1ed9t-cuxe2u-lux1ec7nh-if}

Chương đưa ra một quy trình có hệ thống để xây dựng quyết định.

Giả sử một cửa hàng áp dụng một mức giảm giá cho giá bán dưới một ngưỡng
và mức giảm lớn hơn khi đạt hoặc vượt ngưỡng đó.

\subsection{Bước 1. Xác định điều kiện rẽ
nhánh}\label{bux1b0ux1edbc-1.-xuxe1c-ux111ux1ecbnh-ux111iux1ec1u-kiux1ec7n-rux1ebd-nhuxe1nh}

Đặt một câu hỏi có thể trả lời bằng \texttt{True} hoặc \texttt{False}.

Ví dụ:

\begin{Shaded}
\begin{Highlighting}[]
\NormalTok{Giá ban đầu có nhỏ hơn 128 không?}
\end{Highlighting}
\end{Shaded}

\begin{center}\rule{0.5\linewidth}{0.5pt}\end{center}

\subsection{Bước 2. Mô tả điều xảy ra khi điều kiện
đúng}\label{bux1b0ux1edbc-2.-muxf4-tux1ea3-ux111iux1ec1u-xux1ea3y-ra-khi-ux111iux1ec1u-kiux1ec7n-ux111uxfang}

\begin{Shaded}
\begin{Highlighting}[]
\NormalTok{Dùng mức giảm giá thấp hơn.}
\end{Highlighting}
\end{Shaded}

\begin{center}\rule{0.5\linewidth}{0.5pt}\end{center}

\subsection{Bước 3. Mô tả điều xảy ra khi điều kiện
sai}\label{bux1b0ux1edbc-3.-muxf4-tux1ea3-ux111iux1ec1u-xux1ea3y-ra-khi-ux111iux1ec1u-kiux1ec7n-sai}

\begin{Shaded}
\begin{Highlighting}[]
\NormalTok{Dùng mức giảm giá cao hơn.}
\end{Highlighting}
\end{Shaded}

\begin{center}\rule{0.5\linewidth}{0.5pt}\end{center}

\subsection{Bước 4. Kiểm tra toán tử quan hệ tại giá trị
biên}\label{bux1b0ux1edbc-4.-kiux1ec3m-tra-touxe1n-tux1eed-quan-hux1ec7-tux1ea1i-giuxe1-trux1ecb-biuxean}

Nếu mức giảm cao hơn bắt đầu ngay tại \texttt{128}, điều kiện cho mức
thấp phải là:

\begin{Shaded}
\begin{Highlighting}[]
\NormalTok{originalPrice }\OperatorTok{\textless{}} \DecValTok{128}
\end{Highlighting}
\end{Shaded}

không phải:

\begin{Shaded}
\begin{Highlighting}[]
\NormalTok{originalPrice }\OperatorTok{\textless{}=} \DecValTok{128}
\end{Highlighting}
\end{Shaded}

Các giá trị biên thường giúp phát hiện lỗi chọn \texttt{\textless{}} hay
\texttt{\textless{}=}.

\begin{center}\rule{0.5\linewidth}{0.5pt}\end{center}

\subsection{Bước 5. Loại bỏ phần xử lý lặp
lại}\label{bux1b0ux1edbc-5.-loux1ea1i-bux1ecf-phux1ea7n-xux1eed-luxfd-lux1eb7p-lux1ea1i}

Thay vì tính toàn bộ giá sau giảm trong cả hai nhánh, chỉ chọn hệ số
giảm giá bên trong nhánh:

\begin{Shaded}
\begin{Highlighting}[]
\ControlFlowTok{if}\NormalTok{ originalPrice }\OperatorTok{\textless{}} \DecValTok{128}\NormalTok{ :}
\NormalTok{    discountRate }\OperatorTok{=} \FloatTok{0.92}
\ControlFlowTok{else}\NormalTok{ :}
\NormalTok{    discountRate }\OperatorTok{=} \FloatTok{0.84}

\NormalTok{discountedPrice }\OperatorTok{=}\NormalTok{ discountRate }\OperatorTok{*}\NormalTok{ originalPrice}
\end{Highlighting}
\end{Shaded}

\begin{center}\rule{0.5\linewidth}{0.5pt}\end{center}

\subsection{Bước 6. Kiểm thử cả hai
nhánh}\label{bux1b0ux1edbc-6.-kiux1ec3m-thux1eed-cux1ea3-hai-nhuxe1nh}

Chọn ít nhất một giá trị dưới ngưỡng và một giá trị trên ngưỡng.

Ví dụ:

\begin{Shaded}
\begin{Highlighting}[]
\NormalTok{100  → nhánh giảm thấp hơn}
\NormalTok{200  → nhánh giảm cao hơn}
\NormalTok{128  → giá trị biên}
\end{Highlighting}
\end{Shaded}

\begin{center}\rule{0.5\linewidth}{0.5pt}\end{center}

\subsection{Bước 7. Ghép thành chương trình
Python}\label{bux1b0ux1edbc-7.-ghuxe9p-thuxe0nh-chux1b0ux1a1ng-truxecnh-python}

\begin{Shaded}
\begin{Highlighting}[]
\NormalTok{originalPrice }\OperatorTok{=} \BuiltInTok{float}\NormalTok{(}\BuiltInTok{input}\NormalTok{(}\StringTok{"Original price before discount: "}\NormalTok{))}

\ControlFlowTok{if}\NormalTok{ originalPrice }\OperatorTok{\textless{}} \DecValTok{128}\NormalTok{ :}
\NormalTok{    discountRate }\OperatorTok{=} \FloatTok{0.92}
\ControlFlowTok{else}\NormalTok{ :}
\NormalTok{    discountRate }\OperatorTok{=} \FloatTok{0.84}

\NormalTok{discountedPrice }\OperatorTok{=}\NormalTok{ discountRate }\OperatorTok{*}\NormalTok{ originalPrice}
\BuiltInTok{print}\NormalTok{(}\StringTok{"Discounted price: }\SpecialCharTok{\%.2f}\StringTok{"} \OperatorTok{\%}\NormalTok{ discountedPrice)}
\end{Highlighting}
\end{Shaded}

Điều quan trọng không phải ghi nhớ ví dụ cụ thể này, mà là biết cách
chuyển một quy tắc hai trường hợp thành một điều kiện đúng và hai nhánh
tương ứng.

\begin{center}\rule{0.5\linewidth}{0.5pt}\end{center}

\section{Ví dụ minh họa 3.1 -- Lấy phần giữa của một
chuỗi}\label{vuxed-dux1ee5-minh-hux1ecda-3.1-lux1ea5y-phux1ea7n-giux1eefa-cux1ee7a-mux1ed9t-chuux1ed7i}

Giả sử chương trình cần lấy:

\begin{itemize}
\tightlist
\item
  một ký tự ở giữa nếu độ dài chuỗi là lẻ;
\item
  hai ký tự ở giữa nếu độ dài chuỗi là chẵn.
\end{itemize}

Ví dụ:

\begin{Shaded}
\begin{Highlighting}[]
\NormalTok{"crate"   → "a"}
\NormalTok{"crates"  → "at"}
\end{Highlighting}
\end{Shaded}

Điều kiện rẽ nhánh dựa vào tính chẵn/lẻ:

\begin{Shaded}
\begin{Highlighting}[]
\BuiltInTok{len}\NormalTok{(text) }\OperatorTok{\%} \DecValTok{2} \OperatorTok{==} \DecValTok{1}
\end{Highlighting}
\end{Shaded}

Vị trí giữa có thể tính một lần:

\begin{Shaded}
\begin{Highlighting}[]
\NormalTok{position }\OperatorTok{=} \BuiltInTok{len}\NormalTok{(text) }\OperatorTok{//} \DecValTok{2}
\end{Highlighting}
\end{Shaded}

Sau đó:

\begin{Shaded}
\begin{Highlighting}[]
\ControlFlowTok{if} \BuiltInTok{len}\NormalTok{(text) }\OperatorTok{\%} \DecValTok{2} \OperatorTok{==} \DecValTok{1}\NormalTok{ :}
\NormalTok{    result }\OperatorTok{=}\NormalTok{ text[position]}
\ControlFlowTok{else}\NormalTok{ :}
\NormalTok{    result }\OperatorTok{=}\NormalTok{ text[position }\OperatorTok{{-}} \DecValTok{1}\NormalTok{] }\OperatorTok{+}\NormalTok{ text[position]}
\end{Highlighting}
\end{Shaded}

Ví dụ này kết hợp:

\begin{itemize}
\tightlist
\item
  độ dài chuỗi;
\item
  chia lấy phần nguyên;
\item
  phép chia lấy dư;
\item
  truy cập bằng chỉ số;
\item
  quyết định hai nhánh;
\item
  tránh tính toán lặp lại.
\end{itemize}

\begin{center}\rule{0.5\linewidth}{0.5pt}\end{center}

\subsection{Tự kiểm tra 3.2}\label{tux1ef1-kiux1ec3m-tra-3.2}

\subsubsection{Câu 1}\label{cuxe2u-1-1}

Giá trị của:

\begin{Shaded}
\begin{Highlighting}[]
\DecValTok{4} \OperatorTok{\textless{}=} \DecValTok{4}
\end{Highlighting}
\end{Shaded}

là gì?

\textbf{Đáp án}

\texttt{True}.

\subsubsection{Câu 2}\label{cuxe2u-2-1}

Tại sao đoạn sau sai?

\begin{Shaded}
\begin{Highlighting}[]
\ControlFlowTok{if}\NormalTok{ score }\OperatorTok{=} \DecValTok{100}\NormalTok{ :}
    \BuiltInTok{print}\NormalTok{(}\StringTok{"Perfect"}\NormalTok{)}
\end{Highlighting}
\end{Shaded}

\textbf{Đáp án}

\texttt{=} là phép gán. Muốn kiểm tra bằng nhau phải dùng \texttt{==}.

\subsubsection{Câu 3}\label{cuxe2u-3-1}

Điều kiện nào kiểm tra \texttt{n} là số lẻ?

\textbf{Đáp án}

\begin{Shaded}
\begin{Highlighting}[]
\NormalTok{n }\OperatorTok{\%} \DecValTok{2} \OperatorTok{==} \DecValTok{1}
\end{Highlighting}
\end{Shaded}

\subsubsection{Câu 4}\label{cuxe2u-4-1}

Kết quả phần giữa của:

\begin{Shaded}
\begin{Highlighting}[]
\NormalTok{text }\OperatorTok{=} \StringTok{"monitor"}
\end{Highlighting}
\end{Shaded}

là gì?

\textbf{Đáp án}

\texttt{"i"}.

\begin{center}\rule{0.5\linewidth}{0.5pt}\end{center}

\section{3.3 Rẽ nhánh lồng nhau}\label{rux1ebd-nhuxe1nh-lux1ed3ng-nhau}

Đôi khi chương trình phải đưa ra một quyết định, rồi tiếp tục đưa ra một
quyết định khác bên trong nhánh đã được chọn.

Đó là \textbf{rẽ nhánh lồng nhau}.

Cấu trúc ví dụ:

\begin{Shaded}
\begin{Highlighting}[]
\ControlFlowTok{if}\NormalTok{ condition1 :}
    \ControlFlowTok{if}\NormalTok{ condition2 :}
\NormalTok{        statementA}
    \ControlFlowTok{else}\NormalTok{ :}
\NormalTok{        statementB}
\ControlFlowTok{else}\NormalTok{ :}
\NormalTok{    statementC}
\end{Highlighting}
\end{Shaded}

Câu lệnh \texttt{if} thứ hai nằm bên trong nhánh đầu của \texttt{if} bên
ngoài.

\begin{center}\rule{0.5\linewidth}{0.5pt}\end{center}

\subsection{3.3.1 Quyết định nhiều
tầng}\label{quyux1ebft-ux111ux1ecbnh-nhiux1ec1u-tux1ea7ng}

Một ví dụ điển hình là phép tính thuế đơn giản hóa:

\begin{enumerate}
\def\labelenumi{\arabic{enumi}.}
\tightlist
\item
  trước tiên chương trình xác định nhóm người nộp thuế;
\item
  sau đó xác định mức thu nhập áp dụng trong nhóm đó.
\end{enumerate}

Về mặt ý tưởng:

\begin{Shaded}
\begin{Highlighting}[]
\NormalTok{                     status?}
\NormalTok{                 /             \textbackslash{}}
\NormalTok{             single           married}
\NormalTok{              /                  \textbackslash{}}
\NormalTok{        income limit?        income limit?}
\NormalTok{          /      \textbackslash{}             /      \textbackslash{}}
\NormalTok{      lower     upper       lower     upper}
\NormalTok{      rate      rate        rate      rate}
\end{Highlighting}
\end{Shaded}

Mẫu cài đặt:

\begin{Shaded}
\begin{Highlighting}[]
\ControlFlowTok{if}\NormalTok{ status }\OperatorTok{==} \StringTok{"s"}\NormalTok{ :}
    \ControlFlowTok{if}\NormalTok{ income }\OperatorTok{\textless{}=}\NormalTok{ SINGLE\_LIMIT :}
        \CommentTok{\# mức thuế thấp cho người độc thân}
\NormalTok{        ...}
    \ControlFlowTok{else}\NormalTok{ :}
        \CommentTok{\# mức thuế cao cho người độc thân}
\NormalTok{        ...}
\ControlFlowTok{else}\NormalTok{ :}
    \ControlFlowTok{if}\NormalTok{ income }\OperatorTok{\textless{}=}\NormalTok{ MARRIED\_LIMIT :}
        \CommentTok{\# mức thuế thấp cho người đã kết hôn}
\NormalTok{        ...}
    \ControlFlowTok{else}\NormalTok{ :}
        \CommentTok{\# mức thuế cao cho người đã kết hôn}
\NormalTok{        ...}
\end{Highlighting}
\end{Shaded}

Rẽ nhánh có thể lồng sâu hơn, nhưng quá nhiều mức lồng nhau làm chương
trình khó đọc. Ở các bài sau, hàm sẽ giúp tổ chức logic phức tạp tốt
hơn.

\begin{center}\rule{0.5\linewidth}{0.5pt}\end{center}

\section{Mẹo lập trình 3.2 -- Mô phỏng bằng
tay}\label{mux1eb9o-lux1eadp-truxecnh-3.2-muxf4-phux1ecfng-bux1eb1ng-tay}

\textbf{Mô phỏng bằng tay} (\emph{hand-tracing}) là thực hiện chương
trình thủ công, từng câu lệnh một.

Có thể tạo bảng nhỏ như sau:

\begin{longtable}[]{@{}lrlrr@{}}
\toprule\noalign{}
Bước & \texttt{income} & \texttt{status} & \texttt{tax1} &
\texttt{tax2} \\
\midrule\noalign{}
\endhead
\bottomrule\noalign{}
\endlastfoot
Nhập & 80000 & \texttt{"m"} & & \\
Khởi tạo & 80000 & \texttt{"m"} & 0 & 0 \\
Rẽ nhánh & 80000 & \texttt{"m"} & \ldots{} & \ldots{} \\
\end{longtable}

Mỗi khi một biến thay đổi, cập nhật giá trị của biến trong bảng.

Kỹ thuật này đặc biệt hữu ích khi:

\begin{itemize}
\tightlist
\item
  có nhiều \texttt{if} lồng nhau;
\item
  chương trình có nhiều biến;
\item
  cần hiểu vì sao chương trình đi vào một nhánh nào đó;
\item
  cần so sánh kết quả mong đợi với kết quả thực tế.
\end{itemize}

\begin{center}\rule{0.5\linewidth}{0.5pt}\end{center}

\section{Máy tính và xã hội 3.1 -- Hệ thống phức tạp và logic quyết
định}\label{muxe1y-tuxednh-vuxe0-xuxe3-hux1ed9i-3.1-hux1ec7-thux1ed1ng-phux1ee9c-tux1ea1p-vuxe0-logic-quyux1ebft-ux111ux1ecbnh}

Chương sử dụng hệ thống xử lý hành lý của Sân bay Quốc tế Denver như một
ví dụ về độ phức tạp của phần mềm và hệ thống. Một logic quyết định có
thể rất đơn giản khi xét riêng lẻ, nhưng trở nên khó quản lý khi kết hợp
nhiều thành phần vật lý, ràng buộc thời gian, trường hợp ngoại lệ và
tương tác giữa các thành phần.

Bài học thực tế cho người mới học lập trình là:

\begin{quote}
Mỗi quyết định cần đủ rõ để có thể hiểu và kiểm thử. Không nên giả định
rằng một hệ thống lớn sẽ đúng chỉ vì từng mảnh nhỏ của nó đều có vẻ đúng
khi xét riêng.
\end{quote}

\begin{center}\rule{0.5\linewidth}{0.5pt}\end{center}

\subsection{Tự kiểm tra 3.3}\label{tux1ef1-kiux1ec3m-tra-3.3}

\subsubsection{Câu 1}\label{cuxe2u-1-2}

Khi nào nên dùng rẽ nhánh lồng nhau?

\textbf{Đáp án}

Khi quyết định thứ hai chỉ có ý nghĩa sau khi một kết quả cụ thể của
quyết định trước đã xảy ra.

\subsubsection{Câu 2}\label{cuxe2u-2-2}

Tại sao thụt lề đặc biệt quan trọng trong rẽ nhánh lồng nhau?

\textbf{Đáp án}

Vì thụt lề xác định câu lệnh bên trong thuộc nhánh nào của câu lệnh bên
ngoài.

\subsubsection{Câu 3}\label{cuxe2u-3-2}

Mô phỏng bằng tay giúp hiểu điều gì?

\textbf{Đáp án}

Giúp hiểu thứ tự thực hiện, nhánh nào được chọn và giá trị các biến thay
đổi như thế nào.

\begin{center}\rule{0.5\linewidth}{0.5pt}\end{center}

\section{3.4 Nhiều phương án}\label{nhiux1ec1u-phux1b0ux1a1ng-uxe1n}

Nhiều quyết định có nhiều hơn hai kết quả có thể xảy ra.

Python dùng \texttt{elif} để biểu diễn chuỗi các phương án loại trừ lẫn
nhau.

Cấu trúc tổng quát:

\begin{Shaded}
\begin{Highlighting}[]
\ControlFlowTok{if}\NormalTok{ condition1 :}
\NormalTok{    statements1}
\ControlFlowTok{elif}\NormalTok{ condition2 :}
\NormalTok{    statements2}
\ControlFlowTok{elif}\NormalTok{ condition3 :}
\NormalTok{    statements3}
\ControlFlowTok{else}\NormalTok{ :}
\NormalTok{    defaultStatements}
\end{Highlighting}
\end{Shaded}

Python kiểm tra điều kiện từ trên xuống dưới. Khi gặp điều kiện đầu tiên
đúng, nhánh tương ứng được thực hiện và các nhánh còn lại bị bỏ qua.

\begin{center}\rule{0.5\linewidth}{0.5pt}\end{center}

\subsection{Ví dụ -- Phân loại động
đất}\label{vuxed-dux1ee5-phuxe2n-loux1ea1i-ux111ux1ed9ng-ux111ux1ea5t}

Một cách phân loại đơn giản có thể dùng các ngưỡng:

\begin{Shaded}
\begin{Highlighting}[]
\NormalTok{richter }\OperatorTok{=} \BuiltInTok{float}\NormalTok{(}\BuiltInTok{input}\NormalTok{(}\StringTok{"Enter a magnitude: "}\NormalTok{))}

\ControlFlowTok{if}\NormalTok{ richter }\OperatorTok{\textgreater{}=} \FloatTok{8.0}\NormalTok{ :}
\NormalTok{    description }\OperatorTok{=} \StringTok{"Very severe structural damage"}
\ControlFlowTok{elif}\NormalTok{ richter }\OperatorTok{\textgreater{}=} \FloatTok{7.0}\NormalTok{ :}
\NormalTok{    description }\OperatorTok{=} \StringTok{"Severe damage in many areas"}
\ControlFlowTok{elif}\NormalTok{ richter }\OperatorTok{\textgreater{}=} \FloatTok{6.0}\NormalTok{ :}
\NormalTok{    description }\OperatorTok{=} \StringTok{"Considerable damage is possible"}
\ControlFlowTok{elif}\NormalTok{ richter }\OperatorTok{\textgreater{}=} \FloatTok{4.5}\NormalTok{ :}
\NormalTok{    description }\OperatorTok{=} \StringTok{"Some poorly constructed buildings may be damaged"}
\ControlFlowTok{else}\NormalTok{ :}
\NormalTok{    description }\OperatorTok{=} \StringTok{"Little or no structural damage"}

\BuiltInTok{print}\NormalTok{(description)}
\end{Highlighting}
\end{Shaded}

Điểm cần chú ý là cấu trúc rẽ nhánh, không phải nội dung cụ thể của các
mô tả.

\begin{center}\rule{0.5\linewidth}{0.5pt}\end{center}

\subsection{3.4.1 Thứ tự điều kiện rất quan
trọng}\label{thux1ee9-tux1ef1-ux111iux1ec1u-kiux1ec7n-rux1ea5t-quan-trux1ecdng}

Giả sử điều kiện được viết từ ngưỡng nhỏ lên ngưỡng lớn:

\begin{Shaded}
\begin{Highlighting}[]
\ControlFlowTok{if}\NormalTok{ richter }\OperatorTok{\textgreater{}=} \FloatTok{4.5}\NormalTok{ :}
\NormalTok{    ...}
\ControlFlowTok{elif}\NormalTok{ richter }\OperatorTok{\textgreater{}=} \FloatTok{6.0}\NormalTok{ :}
\NormalTok{    ...}
\ControlFlowTok{elif}\NormalTok{ richter }\OperatorTok{\textgreater{}=} \FloatTok{7.0}\NormalTok{ :}
\NormalTok{    ...}
\end{Highlighting}
\end{Shaded}

Với \texttt{richter\ =\ 7.1}, điều kiện đầu tiên đã đúng nên các nhánh
sau không bao giờ được kiểm tra.

Khi các khoảng chồng lấn, cần kiểm tra \textbf{ngưỡng cao hơn / trường
hợp cụ thể hơn trước}.

Mẫu đúng:

\begin{Shaded}
\begin{Highlighting}[]
\ControlFlowTok{if}\NormalTok{ value }\OperatorTok{\textgreater{}=} \DecValTok{80}\NormalTok{ :}
\NormalTok{    ...}
\ControlFlowTok{elif}\NormalTok{ value }\OperatorTok{\textgreater{}=} \DecValTok{70}\NormalTok{ :}
\NormalTok{    ...}
\ControlFlowTok{elif}\NormalTok{ value }\OperatorTok{\textgreater{}=} \DecValTok{60}\NormalTok{ :}
\NormalTok{    ...}
\ControlFlowTok{else}\NormalTok{ :}
\NormalTok{    ...}
\end{Highlighting}
\end{Shaded}

\begin{center}\rule{0.5\linewidth}{0.5pt}\end{center}

\subsection{\texorpdfstring{3.4.2 \texttt{elif} và các câu lệnh
\texttt{if} độc
lập}{3.4.2 elif và các câu lệnh if độc lập}}\label{elif-vuxe0-cuxe1c-cuxe2u-lux1ec7nh-if-ux111ux1ed9c-lux1eadp}

Hai cách này không tương đương.

Các phương án loại trừ lẫn nhau:

\begin{Shaded}
\begin{Highlighting}[]
\ControlFlowTok{if}\NormalTok{ score }\OperatorTok{\textgreater{}=} \DecValTok{90}\NormalTok{ :}
\NormalTok{    grade }\OperatorTok{=} \StringTok{"A"}
\ControlFlowTok{elif}\NormalTok{ score }\OperatorTok{\textgreater{}=} \DecValTok{80}\NormalTok{ :}
\NormalTok{    grade }\OperatorTok{=} \StringTok{"B"}
\ControlFlowTok{elif}\NormalTok{ score }\OperatorTok{\textgreater{}=} \DecValTok{70}\NormalTok{ :}
\NormalTok{    grade }\OperatorTok{=} \StringTok{"C"}
\end{Highlighting}
\end{Shaded}

Chỉ một nhánh được thực hiện.

Các kiểm tra độc lập:

\begin{Shaded}
\begin{Highlighting}[]
\ControlFlowTok{if}\NormalTok{ score }\OperatorTok{\textgreater{}=} \DecValTok{90}\NormalTok{ :}
    \BuiltInTok{print}\NormalTok{(}\StringTok{"At least 90"}\NormalTok{)}
\ControlFlowTok{if}\NormalTok{ score }\OperatorTok{\textgreater{}=} \DecValTok{80}\NormalTok{ :}
    \BuiltInTok{print}\NormalTok{(}\StringTok{"At least 80"}\NormalTok{)}
\ControlFlowTok{if}\NormalTok{ score }\OperatorTok{\textgreater{}=} \DecValTok{70}\NormalTok{ :}
    \BuiltInTok{print}\NormalTok{(}\StringTok{"At least 70"}\NormalTok{)}
\end{Highlighting}
\end{Shaded}

Với \texttt{score\ =\ 95}, cả ba dòng đều được in.

Dùng \texttt{if/elif/else} khi các trường hợp loại trừ lẫn nhau. Dùng
nhiều \texttt{if} độc lập khi nhiều điều kiện có thể đồng thời kích hoạt
nhiều hành động.

\begin{center}\rule{0.5\linewidth}{0.5pt}\end{center}

\subsection{Tự kiểm tra 3.4}\label{tux1ef1-kiux1ec3m-tra-3.4}

\subsubsection{Câu 1}\label{cuxe2u-1-3}

Viết logic đặt \texttt{sign} bằng \texttt{1}, \texttt{0} hoặc
\texttt{-1} tùy theo \texttt{x} dương, bằng 0 hay âm.

\textbf{Đáp án}

\begin{Shaded}
\begin{Highlighting}[]
\ControlFlowTok{if}\NormalTok{ x }\OperatorTok{\textgreater{}} \DecValTok{0}\NormalTok{ :}
\NormalTok{    sign }\OperatorTok{=} \DecValTok{1}
\ControlFlowTok{elif}\NormalTok{ x }\OperatorTok{\textless{}} \DecValTok{0}\NormalTok{ :}
\NormalTok{    sign }\OperatorTok{=} \OperatorTok{{-}}\DecValTok{1}
\ControlFlowTok{else}\NormalTok{ :}
\NormalTok{    sign }\OperatorTok{=} \DecValTok{0}
\end{Highlighting}
\end{Shaded}

\subsubsection{Câu 2}\label{cuxe2u-2-3}

Tại sao các điều kiện ngưỡng chồng lấn thường nên sắp xếp từ lớn xuống
nhỏ?

\textbf{Đáp án}

Vì Python dừng ở nhánh đúng đầu tiên. Nếu đặt một điều kiện rộng ở
trước, nó có thể chặn một trường hợp cụ thể hơn ở phía sau.

\begin{center}\rule{0.5\linewidth}{0.5pt}\end{center}

\section{TOOLBOX 3.1 -- Gửi e-mail}\label{toolbox-3.1-gux1eedi-e-mail}

Giáo trình dùng tự động hóa e-mail để minh họa cách kết hợp dữ liệu,
quyết định và thư viện Python.

Ý tưởng lập trình cốt lõi là tạo nội dung thông điệp khác nhau tùy theo
điều kiện:

\begin{Shaded}
\begin{Highlighting}[]
\NormalTok{score }\OperatorTok{=} \BuiltInTok{int}\NormalTok{(}\BuiltInTok{input}\NormalTok{(}\StringTok{"Score: "}\NormalTok{))}
\NormalTok{body }\OperatorTok{=} \StringTok{"Your score on the last exam is "} \OperatorTok{+} \BuiltInTok{str}\NormalTok{(score) }\OperatorTok{+} \StringTok{"}\CharTok{\textbackslash{}n}\StringTok{"}

\ControlFlowTok{if}\NormalTok{ score }\OperatorTok{\textless{}=} \DecValTok{50}\NormalTok{ :}
\NormalTok{    body }\OperatorTok{+=} \StringTok{"Please consider visiting the tutoring center."}
\ControlFlowTok{elif}\NormalTok{ score }\OperatorTok{\textgreater{}=} \DecValTok{90}\NormalTok{ :}
\NormalTok{    body }\OperatorTok{+=} \StringTok{"Excellent work."}
\end{Highlighting}
\end{Shaded}

Chi tiết thư viện gửi mail không phải trọng tâm chính của bài. Điều cần
hiểu là chương trình có thể tùy biến hành động hoặc nội dung dựa trên dữ
liệu.

\begin{quote}
Với tài khoản thật, yêu cầu xác thực của nhà cung cấp e-mail có thể thay
đổi theo thời gian. Không nên đặt mật khẩu thật trực tiếp trong mã nguồn
bài giảng hoặc notebook chia sẻ.
\end{quote}

\begin{center}\rule{0.5\linewidth}{0.5pt}\end{center}

\section{3.5 Giải quyết vấn đề: Lưu
đồ}\label{giux1ea3i-quyux1ebft-vux1ea5n-ux111ux1ec1-lux1b0u-ux111ux1ed3}

\textbf{Lưu đồ} (\emph{flowchart}) mô tả trực quan luồng điều khiển của
thuật toán.

Các thành phần thường gặp:

\begin{Shaded}
\begin{Highlighting}[]
\NormalTok{[ Xử lý / công việc ]}

\NormalTok{\textless{} Nhập / xuất \textgreater{}}

\NormalTok{      / Điều kiện? \textbackslash{}}
\NormalTok{     /             \textbackslash{}}
\NormalTok{   Đúng            Sai}
\end{Highlighting}
\end{Shaded}

Lưu đồ hữu ích khi bài toán có nhiều quyết định và chưa rõ cách các
nhánh liên kết với nhau.

\begin{center}\rule{0.5\linewidth}{0.5pt}\end{center}

\subsection{3.5.1 Quyết định hai
nhánh}\label{quyux1ebft-ux111ux1ecbnh-hai-nhuxe1nh-1}

\begin{Shaded}
\begin{Highlighting}[]
\NormalTok{              temperature \textless{} 0 ?}
\NormalTok{                  /        \textbackslash{}}
\NormalTok{               Đúng        Sai}
\NormalTok{                |           |}
\NormalTok{        print("Frozen")    tiếp tục}
\end{Highlighting}
\end{Shaded}

\begin{center}\rule{0.5\linewidth}{0.5pt}\end{center}

\subsection{3.5.2 Nhiều phương
án}\label{nhiux1ec1u-phux1b0ux1a1ng-uxe1n-1}

\begin{Shaded}
\begin{Highlighting}[]
\NormalTok{                score \textgreater{}= 90 ?}
\NormalTok{                 /       \textbackslash{}}
\NormalTok{              Đúng       Sai}
\NormalTok{               |           |}
\NormalTok{             điểm A      score \textgreater{}= 80 ?}
\NormalTok{                         /       \textbackslash{}}
\NormalTok{                      Đúng       Sai}
\NormalTok{                       |           |}
\NormalTok{                    điểm B      ...}
\end{Highlighting}
\end{Shaded}

\begin{center}\rule{0.5\linewidth}{0.5pt}\end{center}

\subsection{3.5.3 Tránh luồng điều khiển kiểu
``spaghetti''}\label{truxe1nh-luux1ed3ng-ux111iux1ec1u-khiux1ec3n-kiux1ec3u-spaghetti}

Một nguyên tắc thiết kế quan trọng là:

\begin{quote}
Không vẽ mũi tên từ một nhánh chui vào giữa một nhánh khác.
\end{quote}

Các nhánh nên được tổ chức rõ ràng: tách ra, thực hiện riêng, rồi có thể
nhập lại sau đó.

Cách tổ chức này phù hợp với \texttt{if}, \texttt{if/else} và các cấu
trúc lồng nhau có cấu trúc rõ ràng.

\begin{center}\rule{0.5\linewidth}{0.5pt}\end{center}

\subsection{Ví dụ -- Chi phí vận
chuyển}\label{vuxed-dux1ee5-chi-phuxed-vux1eadn-chuyux1ec3n}

Giả sử:

\begin{itemize}
\tightlist
\item
  vận chuyển trong phần lục địa của Hoa Kỳ có một mức giá;
\item
  Alaska hoặc Hawaii có mức khác;
\item
  vận chuyển quốc tế dùng mức cao hơn.
\end{itemize}

Một lời giải có cấu trúc:

\begin{Shaded}
\begin{Highlighting}[]
\NormalTok{country }\OperatorTok{=} \BuiltInTok{input}\NormalTok{(}\StringTok{"Enter the country: "}\NormalTok{)}
\NormalTok{state }\OperatorTok{=} \BuiltInTok{input}\NormalTok{(}\StringTok{"Enter the state or province: "}\NormalTok{)}

\ControlFlowTok{if}\NormalTok{ country }\OperatorTok{==} \StringTok{"USA"}\NormalTok{ :}
    \ControlFlowTok{if}\NormalTok{ state }\OperatorTok{==} \StringTok{"AK"} \KeywordTok{or}\NormalTok{ state }\OperatorTok{==} \StringTok{"HI"}\NormalTok{ :}
\NormalTok{        shippingCost }\OperatorTok{=} \FloatTok{10.0}
    \ControlFlowTok{else}\NormalTok{ :}
\NormalTok{        shippingCost }\OperatorTok{=} \FloatTok{5.0}
\ControlFlowTok{else}\NormalTok{ :}
\NormalTok{    shippingCost }\OperatorTok{=} \FloatTok{10.0}

\BuiltInTok{print}\NormalTok{(}\StringTok{"Shipping cost to }\SpecialCharTok{\%s}\StringTok{, }\SpecialCharTok{\%s}\StringTok{: $}\SpecialCharTok{\%.2f}\StringTok{"} \OperatorTok{\%}
\NormalTok{      (state, country, shippingCost))}
\end{Highlighting}
\end{Shaded}

Có thể vẽ lưu đồ trước rồi mới dịch sang Python.

\begin{center}\rule{0.5\linewidth}{0.5pt}\end{center}

\subsection{Tự kiểm tra 3.5}\label{tux1ef1-kiux1ec3m-tra-3.5}

\subsubsection{Câu 1}\label{cuxe2u-1-4}

Mục đích chính của lưu đồ là gì?

\textbf{Đáp án}

Trực quan hóa các công việc, quyết định và các đường đi có thể xảy ra
của chương trình trước hoặc trong khi cài đặt thuật toán.

\subsubsection{Câu 2}\label{cuxe2u-2-4}

Tại sao không nên để mũi tên đi vào giữa một nhánh khác?

\textbf{Đáp án}

Vì nó tạo ra luồng điều khiển thiếu cấu trúc, khó dịch sang chương
trình, khó kiểm thử và khó bảo trì.

\begin{center}\rule{0.5\linewidth}{0.5pt}\end{center}

\section{3.6 Giải quyết vấn đề: Ca kiểm
thử}\label{giux1ea3i-quyux1ebft-vux1ea5n-ux111ux1ec1-ca-kiux1ec3m-thux1eed}

Chương trình có rẽ nhánh phải được kiểm thử bằng các đầu vào đi qua các
đường thực hiện khác nhau.

Chỉ thử một trường hợp thông thường là chưa đủ.

Một kế hoạch kiểm thử tốt hướng tới:

\begin{enumerate}
\def\labelenumi{\arabic{enumi}.}
\tightlist
\item
  \textbf{bao phủ nhánh} -- mỗi nhánh được thực hiện bởi ít nhất một ca
  kiểm thử;
\item
  \textbf{kiểm thử biên} -- kiểm tra các giá trị đúng tại hoặc gần
  ngưỡng quyết định;
\item
  \textbf{kiểm thử đầu vào không hợp lệ} -- khi chương trình có nhiệm vụ
  xác thực dữ liệu.
\end{enumerate}

\begin{center}\rule{0.5\linewidth}{0.5pt}\end{center}

\subsection{3.6.1 Bao phủ nhánh}\label{bao-phux1ee7-nhuxe1nh}

Với quyết định hai nhánh:

\begin{Shaded}
\begin{Highlighting}[]
\ControlFlowTok{if}\NormalTok{ age }\OperatorTok{\textgreater{}=} \DecValTok{18}\NormalTok{ :}
\NormalTok{    category }\OperatorTok{=} \StringTok{"Adult"}
\ControlFlowTok{else}\NormalTok{ :}
\NormalTok{    category }\OperatorTok{=} \StringTok{"Minor"}
\end{Highlighting}
\end{Shaded}

Cần ít nhất một ca cho mỗi kết quả:

\begin{longtable}[]{@{}rll@{}}
\toprule\noalign{}
Đầu vào & Kết quả mong đợi & Mục đích \\
\midrule\noalign{}
\endhead
\bottomrule\noalign{}
\endlastfoot
\texttt{20} & Adult & Nhánh đúng \\
\texttt{15} & Minor & Nhánh sai \\
\end{longtable}

\begin{center}\rule{0.5\linewidth}{0.5pt}\end{center}

\subsection{3.6.2 Các giá trị
biên}\label{cuxe1c-giuxe1-trux1ecb-biuxean}

Ngưỡng ở đây là \texttt{18}.

Các ca hữu ích:

\begin{Shaded}
\begin{Highlighting}[]
\NormalTok{17}
\NormalTok{18}
\NormalTok{19}
\end{Highlighting}
\end{Shaded}

Giá trị \texttt{18} đặc biệt quan trọng vì giúp phân biệt
\texttt{\textgreater{}=} với \texttt{\textgreater{}}.

\begin{center}\rule{0.5\linewidth}{0.5pt}\end{center}

\subsection{3.6.3 Thiết kế ca kiểm thử trước khi lập
trình}\label{thiux1ebft-kux1ebf-ca-kiux1ec3m-thux1eed-trux1b0ux1edbc-khi-lux1eadp-truxecnh}

Viết trước kết quả mong đợi có nhiều lợi ích:

\begin{itemize}
\tightlist
\item
  buộc ta hiểu rõ quy tắc;
\item
  phát hiện điều kiện biên chưa rõ ràng;
\item
  cho phép kiểm chứng chương trình ngay sau khi cài đặt;
\item
  tránh xu hướng chấp nhận mọi kết quả mà chương trình tình cờ in ra.
\end{itemize}

Mẫu bảng kiểm thử:

\begin{longtable}[]{@{}lll@{}}
\toprule\noalign{}
Ca kiểm thử & Kết quả mong đợi & Lý do \\
\midrule\noalign{}
\endhead
\bottomrule\noalign{}
\endlastfoot
trường hợp thường A & \ldots{} & nhánh thứ nhất \\
trường hợp thường B & \ldots{} & nhánh thứ hai \\
giá trị biên & \ldots{} & đúng tại ngưỡng \\
giá trị không hợp lệ & báo lỗi & kiểm tra validation \\
\end{longtable}

\begin{center}\rule{0.5\linewidth}{0.5pt}\end{center}

\section{Mẹo lập trình 3.3 -- Dành thời gian cho kiểm thử và gỡ
lỗi}\label{mux1eb9o-lux1eadp-truxecnh-3.3-duxe0nh-thux1eddi-gian-cho-kiux1ec3m-thux1eed-vuxe0-gux1ee1-lux1ed7i}

Lập trình không chỉ là gõ mã. Một kế hoạch thực tế cần có thời gian cho:

\begin{itemize}
\tightlist
\item
  hiểu và thiết kế lời giải;
\item
  chuẩn bị ca kiểm thử;
\item
  nhập chương trình và sửa lỗi cú pháp;
\item
  chạy thử và gỡ lỗi hành vi không mong đợi.
\end{itemize}

Chương trình càng nhiều nhánh thì càng có nhiều đường thực hiện và
thường cần nhiều công sức kiểm thử hơn.

\begin{center}\rule{0.5\linewidth}{0.5pt}\end{center}

\subsection{Tự kiểm tra 3.6}\label{tux1ef1-kiux1ec3m-tra-3.6}

\subsubsection{Câu 1}\label{cuxe2u-1-5}

Với điều kiện:

\begin{Shaded}
\begin{Highlighting}[]
\ControlFlowTok{if}\NormalTok{ value }\OperatorTok{\textless{}} \DecValTok{100}\NormalTok{ :}
\end{Highlighting}
\end{Shaded}

giá trị biên quan trọng là gì?

\textbf{Đáp án}

\texttt{100}. Các giá trị \texttt{99} và \texttt{101} cũng hữu ích để
kiểm tra lân cận ngưỡng.

\subsubsection{Câu 2}\label{cuxe2u-2-5}

Tại sao một ca kiểm thử thường không đủ cho \texttt{if/else}?

\textbf{Đáp án}

Vì nó có thể chỉ thực hiện một nhánh và để nhánh còn lại chưa từng được
kiểm tra.

\begin{center}\rule{0.5\linewidth}{0.5pt}\end{center}

\section{3.7 Biến Boolean và các toán tử
Boolean}\label{biux1ebfn-boolean-vuxe0-cuxe1c-touxe1n-tux1eed-boolean}

Kiểu Boolean \texttt{bool} chỉ có hai giá trị:

\begin{Shaded}
\begin{Highlighting}[]
\VariableTok{True}
\VariableTok{False}
\end{Highlighting}
\end{Shaded}

Một biến Boolean có thể lưu kết quả của một điều kiện logic.

\begin{Shaded}
\begin{Highlighting}[]
\NormalTok{failed }\OperatorTok{=} \VariableTok{True}

\ControlFlowTok{if}\NormalTok{ failed :}
    \BuiltInTok{print}\NormalTok{(}\StringTok{"The operation failed"}\NormalTok{)}
\end{Highlighting}
\end{Shaded}

Biến Boolean đôi khi được gọi là \textbf{cờ} (\emph{flag}) vì trạng thái
của nó biểu diễn một điều kiện đang đúng hay sai.

\begin{center}\rule{0.5\linewidth}{0.5pt}\end{center}

\subsection{\texorpdfstring{3.7.1 Toán tử
\texttt{and}}{3.7.1 Toán tử and}}\label{touxe1n-tux1eed-and}

\texttt{and} chỉ đúng khi \textbf{cả hai} toán hạng đều đúng.

\begin{Shaded}
\begin{Highlighting}[]
\ControlFlowTok{if}\NormalTok{ temp }\OperatorTok{\textgreater{}} \DecValTok{0} \KeywordTok{and}\NormalTok{ temp }\OperatorTok{\textless{}} \DecValTok{100}\NormalTok{ :}
    \BuiltInTok{print}\NormalTok{(}\StringTok{"Liquid water"}\NormalTok{)}
\end{Highlighting}
\end{Shaded}

Bảng chân trị:

\begin{longtable}[]{@{}lll@{}}
\toprule\noalign{}
A & B & \texttt{A\ and\ B} \\
\midrule\noalign{}
\endhead
\bottomrule\noalign{}
\endlastfoot
\texttt{True} & \texttt{True} & \texttt{True} \\
\texttt{True} & \texttt{False} & \texttt{False} \\
\texttt{False} & \texttt{True} & \texttt{False} \\
\texttt{False} & \texttt{False} & \texttt{False} \\
\end{longtable}

\begin{center}\rule{0.5\linewidth}{0.5pt}\end{center}

\subsection{\texorpdfstring{3.7.2 Toán tử
\texttt{or}}{3.7.2 Toán tử or}}\label{touxe1n-tux1eed-or}

\texttt{or} đúng khi \textbf{ít nhất một} toán hạng đúng.

\begin{Shaded}
\begin{Highlighting}[]
\ControlFlowTok{if}\NormalTok{ temp }\OperatorTok{\textless{}=} \DecValTok{0} \KeywordTok{or}\NormalTok{ temp }\OperatorTok{\textgreater{}=} \DecValTok{100}\NormalTok{ :}
    \BuiltInTok{print}\NormalTok{(}\StringTok{"Not liquid water"}\NormalTok{)}
\end{Highlighting}
\end{Shaded}

Bảng chân trị:

\begin{longtable}[]{@{}lll@{}}
\toprule\noalign{}
A & B & \texttt{A\ or\ B} \\
\midrule\noalign{}
\endhead
\bottomrule\noalign{}
\endlastfoot
\texttt{True} & \texttt{True} & \texttt{True} \\
\texttt{True} & \texttt{False} & \texttt{True} \\
\texttt{False} & \texttt{True} & \texttt{True} \\
\texttt{False} & \texttt{False} & \texttt{False} \\
\end{longtable}

Lưu ý: \texttt{or} là phép ``hoặc bao hàm'', tức là vẫn đúng khi
\textbf{cả hai} điều kiện đều đúng.

\begin{center}\rule{0.5\linewidth}{0.5pt}\end{center}

\subsection{\texorpdfstring{3.7.3 Toán tử
\texttt{not}}{3.7.3 Toán tử not}}\label{touxe1n-tux1eed-not}

\texttt{not} đảo giá trị Boolean.

\begin{Shaded}
\begin{Highlighting}[]
\ControlFlowTok{if} \KeywordTok{not}\NormalTok{ frozen :}
    \BuiltInTok{print}\NormalTok{(}\StringTok{"Not frozen"}\NormalTok{)}
\end{Highlighting}
\end{Shaded}

\begin{longtable}[]{@{}ll@{}}
\toprule\noalign{}
A & \texttt{not\ A} \\
\midrule\noalign{}
\endhead
\bottomrule\noalign{}
\endlastfoot
\texttt{True} & \texttt{False} \\
\texttt{False} & \texttt{True} \\
\end{longtable}

\begin{center}\rule{0.5\linewidth}{0.5pt}\end{center}

\subsection{3.7.4 Thứ tự ưu tiên}\label{thux1ee9-tux1ef1-ux1b0u-tiuxean}

Một thứ tự đơn giản để ghi nhớ:

\begin{Shaded}
\begin{Highlighting}[]
\NormalTok{số học}
\NormalTok{  ↓}
\NormalTok{toán tử quan hệ}
\NormalTok{  ↓}
\NormalTok{not}
\NormalTok{  ↓}
\NormalTok{and}
\NormalTok{  ↓}
\NormalTok{or}
\end{Highlighting}
\end{Shaded}

Do đó:

\begin{Shaded}
\begin{Highlighting}[]
\NormalTok{x }\OperatorTok{\textgreater{}} \DecValTok{0} \KeywordTok{and}\NormalTok{ y }\OperatorTok{\textless{}} \DecValTok{10}
\end{Highlighting}
\end{Shaded}

được hiểu là:

\begin{Shaded}
\begin{Highlighting}[]
\NormalTok{(x }\OperatorTok{\textgreater{}} \DecValTok{0}\NormalTok{) }\KeywordTok{and}\NormalTok{ (y }\OperatorTok{\textless{}} \DecValTok{10}\NormalTok{)}
\end{Highlighting}
\end{Shaded}

Với biểu thức logic phức tạp, dấu ngoặc vẫn nên được dùng để tăng tính
dễ đọc.

\begin{center}\rule{0.5\linewidth}{0.5pt}\end{center}

\section{\texorpdfstring{Lỗi thường gặp 3.3 -- Nhầm \texttt{and} với
\texttt{or}}{Lỗi thường gặp 3.3 -- Nhầm and với or}}\label{lux1ed7i-thux1b0ux1eddng-gux1eb7p-3.3-nhux1ea7m-and-vux1edbi-or}

Giả sử một giá trị hợp lệ khi nằm trong khoảng từ \texttt{0} đến
\texttt{100}.

Đúng:

\begin{Shaded}
\begin{Highlighting}[]
\ControlFlowTok{if}\NormalTok{ value }\OperatorTok{\textgreater{}=} \DecValTok{0} \KeywordTok{and}\NormalTok{ value }\OperatorTok{\textless{}=} \DecValTok{100}\NormalTok{ :}
    \BuiltInTok{print}\NormalTok{(}\StringTok{"Valid"}\NormalTok{)}
\end{Highlighting}
\end{Shaded}

Cả hai điều kiện phải cùng đúng.

Để phát hiện giá trị không hợp lệ:

\begin{Shaded}
\begin{Highlighting}[]
\ControlFlowTok{if}\NormalTok{ value }\OperatorTok{\textless{}} \DecValTok{0} \KeywordTok{or}\NormalTok{ value }\OperatorTok{\textgreater{}} \DecValTok{100}\NormalTok{ :}
    \BuiltInTok{print}\NormalTok{(}\StringTok{"Invalid"}\NormalTok{)}
\end{Highlighting}
\end{Shaded}

Chỉ cần một trong các điều kiện không hợp lệ là đúng.

Câu hỏi gợi ý:

\begin{itemize}
\tightlist
\item
  Có cần \textbf{tất cả} điều kiện cùng đúng không? → dùng \texttt{and}.
\item
  Chỉ cần \textbf{một} điều kiện đúng là đủ không? → dùng \texttt{or}.
\end{itemize}

\begin{center}\rule{0.5\linewidth}{0.5pt}\end{center}

\section{Mẹo lập trình 3.4 -- Tính dễ
đọc}\label{mux1eb9o-lux1eadp-truxecnh-3.4-tuxednh-dux1ec5-ux111ux1ecdc}

Không nên so sánh biến Boolean một cách thừa thãi với \texttt{True} hoặc
\texttt{False}.

Kém rõ ràng:

\begin{Shaded}
\begin{Highlighting}[]
\ControlFlowTok{if}\NormalTok{ frozen }\OperatorTok{==} \VariableTok{False}\NormalTok{ :}
    \BuiltInTok{print}\NormalTok{(}\StringTok{"Not frozen"}\NormalTok{)}
\end{Highlighting}
\end{Shaded}

Nên viết:

\begin{Shaded}
\begin{Highlighting}[]
\ControlFlowTok{if} \KeywordTok{not}\NormalTok{ frozen :}
    \BuiltInTok{print}\NormalTok{(}\StringTok{"Not frozen"}\NormalTok{)}
\end{Highlighting}
\end{Shaded}

Tương tự:

\begin{Shaded}
\begin{Highlighting}[]
\ControlFlowTok{if}\NormalTok{ valid :}
    \BuiltInTok{print}\NormalTok{(}\StringTok{"Input accepted"}\NormalTok{)}
\end{Highlighting}
\end{Shaded}

thường dễ đọc hơn:

\begin{Shaded}
\begin{Highlighting}[]
\ControlFlowTok{if}\NormalTok{ valid }\OperatorTok{==} \VariableTok{True}\NormalTok{ :}
    \BuiltInTok{print}\NormalTok{(}\StringTok{"Input accepted"}\NormalTok{)}
\end{Highlighting}
\end{Shaded}

Tên biến Boolean nên mô tả rõ một trạng thái hoặc điều kiện:

\begin{Shaded}
\begin{Highlighting}[]
\NormalTok{valid}
\NormalTok{finished}
\NormalTok{found}
\NormalTok{isEmpty}
\NormalTok{hasError}
\end{Highlighting}
\end{Shaded}

\begin{center}\rule{0.5\linewidth}{0.5pt}\end{center}

\section{Chủ đề bổ sung 3.3 -- Nối chuỗi các phép so
sánh}\label{chux1ee7-ux111ux1ec1-bux1ed5-sung-3.3-nux1ed1i-chuux1ed7i-cuxe1c-phuxe9p-so-suxe1nh}

Python cho phép viết:

\begin{Shaded}
\begin{Highlighting}[]
\DecValTok{0} \OperatorTok{\textless{}=}\NormalTok{ value }\OperatorTok{\textless{}=} \DecValTok{100}
\end{Highlighting}
\end{Shaded}

Tương đương với:

\begin{Shaded}
\begin{Highlighting}[]
\NormalTok{value }\OperatorTok{\textgreater{}=} \DecValTok{0} \KeywordTok{and}\NormalTok{ value }\OperatorTok{\textless{}=} \DecValTok{100}
\end{Highlighting}
\end{Shaded}

Cách nối này gọn và tự nhiên trong Python, nhưng người mới học nên hiểu
rõ biểu thức Boolean tương đương.

\begin{center}\rule{0.5\linewidth}{0.5pt}\end{center}

\section{Chủ đề bổ sung 3.4 -- Đánh giá ngắn
mạch}\label{chux1ee7-ux111ux1ec1-bux1ed5-sung-3.4-ux111uxe1nh-giuxe1-ngux1eafn-mux1ea1ch}

Python đánh giá \texttt{and} và \texttt{or} từ trái sang phải và dừng
ngay khi đã xác định được kết quả cuối cùng.

Ví dụ:

\begin{Shaded}
\begin{Highlighting}[]
\NormalTok{quantity }\OperatorTok{\textgreater{}} \DecValTok{0} \KeywordTok{and}\NormalTok{ price }\OperatorTok{/}\NormalTok{ quantity }\OperatorTok{\textless{}} \DecValTok{10}
\end{Highlighting}
\end{Shaded}

Nếu \texttt{quantity\ \textgreater{}\ 0} là sai, Python không thực hiện
phép chia. Điều này tránh chia cho 0 khi \texttt{quantity\ ==\ 0}.

Với \texttt{or}:

\begin{Shaded}
\begin{Highlighting}[]
\NormalTok{conditionA }\KeywordTok{or}\NormalTok{ conditionB}
\end{Highlighting}
\end{Shaded}

nếu \texttt{conditionA} đã đúng, \texttt{conditionB} sẽ không được đánh
giá.

Cơ chế này gọi là \textbf{đánh giá ngắn mạch} (\emph{short-circuit
evaluation}).

\begin{center}\rule{0.5\linewidth}{0.5pt}\end{center}

\section{Chủ đề bổ sung 3.5 -- Định luật De
Morgan}\label{chux1ee7-ux111ux1ec1-bux1ed5-sung-3.5-ux111ux1ecbnh-luux1eadt-de-morgan}

Định luật De Morgan giúp biến đổi các điều kiện ghép có phủ định.

\begin{Shaded}
\begin{Highlighting}[]
\NormalTok{not (A and B)  ≡  (not A) or (not B)}

\NormalTok{not (A or B)   ≡  (not A) and (not B)}
\end{Highlighting}
\end{Shaded}

Ví dụ:

\begin{Shaded}
\begin{Highlighting}[]
\KeywordTok{not}\NormalTok{ (state }\OperatorTok{==} \StringTok{"AK"} \KeywordTok{or}\NormalTok{ state }\OperatorTok{==} \StringTok{"HI"}\NormalTok{)}
\end{Highlighting}
\end{Shaded}

tương đương với:

\begin{Shaded}
\begin{Highlighting}[]
\NormalTok{state }\OperatorTok{!=} \StringTok{"AK"} \KeywordTok{and}\NormalTok{ state }\OperatorTok{!=} \StringTok{"HI"}
\end{Highlighting}
\end{Shaded}

Ví dụ khác:

\begin{Shaded}
\begin{Highlighting}[]
\KeywordTok{not}\NormalTok{ (country }\OperatorTok{==} \StringTok{"USA"} \KeywordTok{and}\NormalTok{ state }\OperatorTok{!=} \StringTok{"AK"} \KeywordTok{and}\NormalTok{ state }\OperatorTok{!=} \StringTok{"HI"}\NormalTok{)}
\end{Highlighting}
\end{Shaded}

có thể viết lại trực tiếp hơn:

\begin{Shaded}
\begin{Highlighting}[]
\NormalTok{country }\OperatorTok{!=} \StringTok{"USA"} \KeywordTok{or}\NormalTok{ state }\OperatorTok{==} \StringTok{"AK"} \KeywordTok{or}\NormalTok{ state }\OperatorTok{==} \StringTok{"HI"}
\end{Highlighting}
\end{Shaded}

Mục tiêu không chỉ là viết ngắn hơn mà là làm logic rõ ràng hơn.

\begin{center}\rule{0.5\linewidth}{0.5pt}\end{center}

\subsection{Tự kiểm tra 3.7}\label{tux1ef1-kiux1ec3m-tra-3.7}

\subsubsection{Câu 1}\label{cuxe2u-1-6}

Làm thế nào kiểm tra cả \texttt{x} và \texttt{y} đều dương?

\textbf{Đáp án}

\begin{Shaded}
\begin{Highlighting}[]
\NormalTok{x }\OperatorTok{\textgreater{}} \DecValTok{0} \KeywordTok{and}\NormalTok{ y }\OperatorTok{\textgreater{}} \DecValTok{0}
\end{Highlighting}
\end{Shaded}

\subsubsection{Câu 2}\label{cuxe2u-2-6}

Làm thế nào kiểm tra ít nhất một trong \texttt{x}, \texttt{y} bằng 0?

\textbf{Đáp án}

\begin{Shaded}
\begin{Highlighting}[]
\NormalTok{x }\OperatorTok{==} \DecValTok{0} \KeywordTok{or}\NormalTok{ y }\OperatorTok{==} \DecValTok{0}
\end{Highlighting}
\end{Shaded}

\subsubsection{Câu 3}\label{cuxe2u-3-3}

Giá trị của:

\begin{Shaded}
\begin{Highlighting}[]
\KeywordTok{not} \KeywordTok{not} \VariableTok{True}
\end{Highlighting}
\end{Shaded}

là gì?

\textbf{Đáp án}

\texttt{True}.

\subsubsection{Câu 4}\label{cuxe2u-4-2}

Tại sao biểu thức sau vẫn an toàn khi \texttt{quantity\ ==\ 0}?

\begin{Shaded}
\begin{Highlighting}[]
\NormalTok{quantity }\OperatorTok{\textgreater{}} \DecValTok{0} \KeywordTok{and}\NormalTok{ price }\OperatorTok{/}\NormalTok{ quantity }\OperatorTok{\textless{}} \DecValTok{10}
\end{Highlighting}
\end{Shaded}

\textbf{Đáp án}

Vì đánh giá ngắn mạch dừng ngay sau điều kiện đầu tiên sai, nên phép
chia không được thực hiện.

\begin{center}\rule{0.5\linewidth}{0.5pt}\end{center}

\section{3.8 Phân tích chuỗi}\label{phuxe2n-tuxedch-chuux1ed7i}

Logic rẽ nhánh thường được dùng để kiểm tra dữ liệu văn bản.

Python cung cấp nhiều toán tử và phương thức để xác định một chuỗi có
chứa chuỗi con hay có đặc điểm nhất định hay không.

\begin{center}\rule{0.5\linewidth}{0.5pt}\end{center}

\subsection{\texorpdfstring{3.8.1 Toán tử thành viên \texttt{in} và
\texttt{not\ in}}{3.8.1 Toán tử thành viên in và not in}}\label{touxe1n-tux1eed-thuxe0nh-viuxean-in-vuxe0-not-in}

\begin{Shaded}
\begin{Highlighting}[]
\NormalTok{name }\OperatorTok{=} \StringTok{"John Wayne"}

\ControlFlowTok{if} \StringTok{"Way"} \KeywordTok{in}\NormalTok{ name :}
    \BuiltInTok{print}\NormalTok{(}\StringTok{"Substring found"}\NormalTok{)}
\end{Highlighting}
\end{Shaded}

Kiểm tra ngược lại:

\begin{Shaded}
\begin{Highlighting}[]
\ControlFlowTok{if} \StringTok{"{-}"} \KeywordTok{not} \KeywordTok{in}\NormalTok{ name :}
    \BuiltInTok{print}\NormalTok{(}\StringTok{"The name does not contain a hyphen"}\NormalTok{)}
\end{Highlighting}
\end{Shaded}

Kiểm tra thành viên phân biệt chữ hoa/chữ thường.

\begin{center}\rule{0.5\linewidth}{0.5pt}\end{center}

\subsection{3.8.2 Một số phương thức chuỗi con hữu
ích}\label{mux1ed9t-sux1ed1-phux1b0ux1a1ng-thux1ee9c-chuux1ed7i-con-hux1eefu-uxedch}

\begin{longtable}[]{@{}
  >{\raggedright\arraybackslash}p{(\columnwidth - 2\tabcolsep) * \real{0.5000}}
  >{\raggedright\arraybackslash}p{(\columnwidth - 2\tabcolsep) * \real{0.5000}}@{}}
\toprule\noalign{}
\begin{minipage}[b]{\linewidth}\raggedright
Phép toán
\end{minipage} & \begin{minipage}[b]{\linewidth}\raggedright
Ý nghĩa
\end{minipage} \\
\midrule\noalign{}
\endhead
\bottomrule\noalign{}
\endlastfoot
\texttt{substring\ in\ s} & \texttt{s} có chứa \texttt{substring} hay
không? \\
\texttt{s.count(substring)} & Số lần xuất hiện không chồng lấn \\
\texttt{s.startswith(substring)} & \texttt{s} có bắt đầu bằng chuỗi con
không? \\
\texttt{s.endswith(substring)} & \texttt{s} có kết thúc bằng chuỗi con
không? \\
\texttt{s.find(substring)} & Chỉ số lần xuất hiện đầu tiên, hoặc
\texttt{-1} nếu không tìm thấy \\
\end{longtable}

Ví dụ:

\begin{Shaded}
\begin{Highlighting}[]
\NormalTok{filename }\OperatorTok{=} \StringTok{"report.html"}

\ControlFlowTok{if}\NormalTok{ filename.endswith(}\StringTok{".html"}\NormalTok{) :}
    \BuiltInTok{print}\NormalTok{(}\StringTok{"HTML file"}\NormalTok{)}
\end{Highlighting}
\end{Shaded}

Ví dụ:

\begin{Shaded}
\begin{Highlighting}[]
\NormalTok{text }\OperatorTok{=} \StringTok{"banana"}
\BuiltInTok{print}\NormalTok{(text.count(}\StringTok{"an"}\NormalTok{))}
\BuiltInTok{print}\NormalTok{(text.find(}\StringTok{"na"}\NormalTok{))}
\end{Highlighting}
\end{Shaded}

\begin{center}\rule{0.5\linewidth}{0.5pt}\end{center}

\subsection{3.8.3 Kiểm tra đặc điểm của
chuỗi}\label{kiux1ec3m-tra-ux111ux1eb7c-ux111iux1ec3m-cux1ee7a-chuux1ed7i}

Một số phương thức hữu ích:

\begin{longtable}[]{@{}
  >{\raggedright\arraybackslash}p{(\columnwidth - 2\tabcolsep) * \real{0.5000}}
  >{\raggedright\arraybackslash}p{(\columnwidth - 2\tabcolsep) * \real{0.5000}}@{}}
\toprule\noalign{}
\begin{minipage}[b]{\linewidth}\raggedright
Phương thức
\end{minipage} & \begin{minipage}[b]{\linewidth}\raggedright
Trả về \texttt{True} khi\ldots{}
\end{minipage} \\
\midrule\noalign{}
\endhead
\bottomrule\noalign{}
\endlastfoot
\texttt{s.isalnum()} & mọi ký tự đều là chữ hoặc số và chuỗi không
rỗng \\
\texttt{s.isalpha()} & mọi ký tự đều là chữ cái và chuỗi không rỗng \\
\texttt{s.isdigit()} & mọi ký tự đều là chữ số và chuỗi không rỗng \\
\texttt{s.islower()} & có ít nhất một chữ cái và mọi chữ cái đều viết
thường \\
\texttt{s.isupper()} & có ít nhất một chữ cái và mọi chữ cái đều viết
hoa \\
\texttt{s.isspace()} & mọi ký tự đều là khoảng trắng và chuỗi không
rỗng \\
\end{longtable}

Ví dụ:

\begin{Shaded}
\begin{Highlighting}[]
\CommentTok{"1729"}\NormalTok{.isdigit()}
\end{Highlighting}
\end{Shaded}

trả về \texttt{True}.

\begin{Shaded}
\begin{Highlighting}[]
\CommentTok{"{-}1729"}\NormalTok{.isdigit()}
\end{Highlighting}
\end{Shaded}

trả về \texttt{False} vì \texttt{-} không phải chữ số.

\begin{Shaded}
\begin{Highlighting}[]
\CommentTok{"John Smith"}\NormalTok{.isalpha()}
\end{Highlighting}
\end{Shaded}

trả về \texttt{False} vì chuỗi chứa khoảng trắng.

\begin{center}\rule{0.5\linewidth}{0.5pt}\end{center}

\subsection{Ví dụ -- Khảo sát chuỗi
con}\label{vuxed-dux1ee5-khux1ea3o-suxe1t-chuux1ed7i-con}

\begin{Shaded}
\begin{Highlighting}[]
\NormalTok{theString }\OperatorTok{=} \BuiltInTok{input}\NormalTok{(}\StringTok{"Enter a string: "}\NormalTok{)}
\NormalTok{theSubstring }\OperatorTok{=} \BuiltInTok{input}\NormalTok{(}\StringTok{"Enter a substring: "}\NormalTok{)}

\ControlFlowTok{if}\NormalTok{ theSubstring }\KeywordTok{in}\NormalTok{ theString :}
    \BuiltInTok{print}\NormalTok{(}\StringTok{"The string contains the substring."}\NormalTok{)}
    \BuiltInTok{print}\NormalTok{(}\StringTok{"Count:"}\NormalTok{, theString.count(theSubstring))}
    \BuiltInTok{print}\NormalTok{(}\StringTok{"First position:"}\NormalTok{, theString.find(theSubstring))}

    \ControlFlowTok{if}\NormalTok{ theString.startswith(theSubstring) :}
        \BuiltInTok{print}\NormalTok{(}\StringTok{"It appears at the beginning."}\NormalTok{)}

    \ControlFlowTok{if}\NormalTok{ theString.endswith(theSubstring) :}
        \BuiltInTok{print}\NormalTok{(}\StringTok{"It appears at the end."}\NormalTok{)}
\ControlFlowTok{else}\NormalTok{ :}
    \BuiltInTok{print}\NormalTok{(}\StringTok{"The string does not contain the substring."}\NormalTok{)}
\end{Highlighting}
\end{Shaded}

\begin{center}\rule{0.5\linewidth}{0.5pt}\end{center}

\subsection{Tự kiểm tra 3.8}\label{tux1ef1-kiux1ec3m-tra-3.8}

\subsubsection{Câu 1}\label{cuxe2u-1-7}

Làm thế nào đếm số khoảng trắng trong \texttt{text}?

\textbf{Đáp án}

\begin{Shaded}
\begin{Highlighting}[]
\NormalTok{text.count(}\StringTok{" "}\NormalTok{)}
\end{Highlighting}
\end{Shaded}

\subsubsection{Câu 2}\label{cuxe2u-2-7}

Làm thế nào kiểm tra tên file kết thúc bằng \texttt{.jpg} hoặc
\texttt{.jpeg}?

\textbf{Đáp án}

\begin{Shaded}
\begin{Highlighting}[]
\NormalTok{filename.endswith(}\StringTok{".jpg"}\NormalTok{) }\KeywordTok{or}\NormalTok{ filename.endswith(}\StringTok{".jpeg"}\NormalTok{)}
\end{Highlighting}
\end{Shaded}

Nếu muốn không phân biệt chữ hoa/thường, hãy chuẩn hóa bằng
\texttt{lower()} trước.

\subsubsection{Câu 3}\label{cuxe2u-3-4}

\texttt{find()} trả về gì khi không tìm thấy chuỗi con?

\textbf{Đáp án}

\texttt{-1}.

\begin{center}\rule{0.5\linewidth}{0.5pt}\end{center}

\section{3.9 Ứng dụng: Kiểm tra dữ liệu đầu
vào}\label{ux1ee9ng-dux1ee5ng-kiux1ec3m-tra-dux1eef-liux1ec7u-ux111ux1ea7u-vuxe0o}

\textbf{Kiểm tra dữ liệu đầu vào} (\emph{input validation}) là xác minh
dữ liệu do người dùng cung cấp trước khi sử dụng.

Chương trình không nên giả định mọi dữ liệu nhập đều thỏa mãn yêu cầu.

Giả sử thang máy chấp nhận các tầng từ \texttt{1} đến \texttt{20}, trừ
tầng \texttt{13}.

Các đầu vào không hợp lệ gồm:

\begin{itemize}
\tightlist
\item
  \texttt{13};
\item
  \texttt{0} hoặc số âm;
\item
  số lớn hơn \texttt{20}.
\end{itemize}

Phiên bản có validation:

\begin{Shaded}
\begin{Highlighting}[]
\NormalTok{floor }\OperatorTok{=} \BuiltInTok{int}\NormalTok{(}\BuiltInTok{input}\NormalTok{(}\StringTok{"Floor: "}\NormalTok{))}

\ControlFlowTok{if}\NormalTok{ floor }\OperatorTok{==} \DecValTok{13}\NormalTok{ :}
    \BuiltInTok{print}\NormalTok{(}\StringTok{"Error: There is no thirteenth floor."}\NormalTok{)}
\ControlFlowTok{elif}\NormalTok{ floor }\OperatorTok{\textless{}=} \DecValTok{0} \KeywordTok{or}\NormalTok{ floor }\OperatorTok{\textgreater{}} \DecValTok{20}\NormalTok{ :}
    \BuiltInTok{print}\NormalTok{(}\StringTok{"Error: The floor must be between 1 and 20."}\NormalTok{)}
\ControlFlowTok{else}\NormalTok{ :}
\NormalTok{    actualFloor }\OperatorTok{=}\NormalTok{ floor}

    \ControlFlowTok{if}\NormalTok{ floor }\OperatorTok{\textgreater{}} \DecValTok{13}\NormalTok{ :}
\NormalTok{        actualFloor }\OperatorTok{=}\NormalTok{ floor }\OperatorTok{{-}} \DecValTok{1}

    \BuiltInTok{print}\NormalTok{(}\StringTok{"The elevator will travel to the actual floor"}\NormalTok{, actualFloor)}
\end{Highlighting}
\end{Shaded}

Mẫu tổng quát:

\begin{Shaded}
\begin{Highlighting}[]
\NormalTok{Đọc dữ liệu}
\NormalTok{   ↓}
\NormalTok{Kiểm tra hợp lệ}
\NormalTok{   ↓}
\NormalTok{Không hợp lệ? → báo lỗi}
\NormalTok{Hợp lệ?       → xử lý}
\end{Highlighting}
\end{Shaded}

\begin{center}\rule{0.5\linewidth}{0.5pt}\end{center}

\subsection{3.9.1 Kiểm tra miền giá
trị}\label{kiux1ec3m-tra-miux1ec1n-giuxe1-trux1ecb}

Ví dụ:

\begin{Shaded}
\begin{Highlighting}[]
\ControlFlowTok{if}\NormalTok{ age }\OperatorTok{\textless{}} \DecValTok{0} \KeywordTok{or}\NormalTok{ age }\OperatorTok{\textgreater{}} \DecValTok{130}\NormalTok{ :}
    \BuiltInTok{print}\NormalTok{(}\StringTok{"Invalid age"}\NormalTok{)}
\ControlFlowTok{else}\NormalTok{ :}
    \CommentTok{\# Có thể xử lý giá trị an toàn hơn.}
\NormalTok{    ...}
\end{Highlighting}
\end{Shaded}

\begin{center}\rule{0.5\linewidth}{0.5pt}\end{center}

\subsection{3.9.2 Kiểm tra lựa chọn nhập
vào}\label{kiux1ec3m-tra-lux1ef1a-chux1ecdn-nhux1eadp-vuxe0o}

Giả sử người dùng phải nhập \texttt{s} hoặc \texttt{m}.

Một cách:

\begin{Shaded}
\begin{Highlighting}[]
\NormalTok{status }\OperatorTok{=} \BuiltInTok{input}\NormalTok{(}\StringTok{"Enter s or m: "}\NormalTok{)}

\ControlFlowTok{if}\NormalTok{ status }\OperatorTok{==} \StringTok{"s"} \KeywordTok{or}\NormalTok{ status }\OperatorTok{==} \StringTok{"S"}\NormalTok{ :}
\NormalTok{    ...}
\ControlFlowTok{elif}\NormalTok{ status }\OperatorTok{==} \StringTok{"m"} \KeywordTok{or}\NormalTok{ status }\OperatorTok{==} \StringTok{"M"}\NormalTok{ :}
\NormalTok{    ...}
\ControlFlowTok{else}\NormalTok{ :}
    \BuiltInTok{print}\NormalTok{(}\StringTok{"Invalid status"}\NormalTok{)}
\end{Highlighting}
\end{Shaded}

Cách dễ mở rộng hơn là chuẩn hóa trước:

\begin{Shaded}
\begin{Highlighting}[]
\NormalTok{status }\OperatorTok{=} \BuiltInTok{input}\NormalTok{(}\StringTok{"Enter s or m: "}\NormalTok{)}
\NormalTok{status }\OperatorTok{=}\NormalTok{ status.lower()}

\ControlFlowTok{if}\NormalTok{ status }\OperatorTok{==} \StringTok{"s"}\NormalTok{ :}
\NormalTok{    ...}
\ControlFlowTok{elif}\NormalTok{ status }\OperatorTok{==} \StringTok{"m"}\NormalTok{ :}
\NormalTok{    ...}
\ControlFlowTok{else}\NormalTok{ :}
    \BuiltInTok{print}\NormalTok{(}\StringTok{"Invalid status"}\NormalTok{)}
\end{Highlighting}
\end{Shaded}

Tương tự, mã nhiều ký tự có thể chuẩn hóa bằng:

\begin{Shaded}
\begin{Highlighting}[]
\NormalTok{country }\OperatorTok{=} \BuiltInTok{input}\NormalTok{(}\StringTok{"Country code: "}\NormalTok{)}
\NormalTok{country }\OperatorTok{=}\NormalTok{ country.upper()}
\end{Highlighting}
\end{Shaded}

Sau đó chỉ cần so sánh với một dạng chuẩn.

\begin{center}\rule{0.5\linewidth}{0.5pt}\end{center}

\subsection{3.9.3 Giới hạn ở giai đoạn hiện
tại}\label{giux1edbi-hux1ea1n-ux1edf-giai-ux111oux1ea1n-hiux1ec7n-tux1ea1i}

Đoạn mã:

\begin{Shaded}
\begin{Highlighting}[]
\NormalTok{floor }\OperatorTok{=} \BuiltInTok{int}\NormalTok{(}\BuiltInTok{input}\NormalTok{(}\StringTok{"Floor: "}\NormalTok{))}
\end{Highlighting}
\end{Shaded}

vẫn phát sinh ngoại lệ nếu người dùng nhập:

\begin{Shaded}
\begin{Highlighting}[]
\NormalTok{five}
\end{Highlighting}
\end{Shaded}

Xử lý lỗi chuyển đổi kiểu cần cơ chế exception, được giới thiệu ở chương
sau.

Ở giai đoạn này cần phân biệt:

\begin{itemize}
\tightlist
\item
  \textbf{kiểm tra giá trị sau khi chuyển đổi} -- ví dụ kiểm tra có nằm
  trong khoảng hợp lệ không;
\item
  \textbf{xử lý việc chuyển đổi thất bại} -- chủ đề về ngoại lệ học ở
  phần sau.
\end{itemize}

\begin{center}\rule{0.5\linewidth}{0.5pt}\end{center}

\section{Chủ đề bổ sung 3.6 -- Kết thúc chương
trình}\label{chux1ee7-ux111ux1ec1-bux1ed5-sung-3.6-kux1ebft-thuxfac-chux1b0ux1a1ng-truxecnh}

Với chương trình văn bản nhỏ, dữ liệu sai đôi khi có thể xử lý bằng cách
kết thúc chương trình ngay.

\begin{Shaded}
\begin{Highlighting}[]
\ImportTok{from}\NormalTok{ sys }\ImportTok{import}\NormalTok{ exit}

\NormalTok{response }\OperatorTok{=} \BuiltInTok{input}\NormalTok{(}\StringTok{"Enter y or n: "}\NormalTok{)}

\ControlFlowTok{if} \KeywordTok{not}\NormalTok{ (response }\OperatorTok{==} \StringTok{"y"} \KeywordTok{or}\NormalTok{ response }\OperatorTok{==} \StringTok{"n"}\NormalTok{) :}
\NormalTok{    exit(}\StringTok{"Error: you must enter y or n."}\NormalTok{)}
\end{Highlighting}
\end{Shaded}

Sau khi validation thành công, phần mã phía sau có thể giả định đầu vào
đã thỏa điều kiện.

Không nên lạm dụng việc kết thúc sớm. Trong chương trình lớn hơn, hàm và
ngoại lệ thường giúp tổ chức mã tốt hơn.

\begin{center}\rule{0.5\linewidth}{0.5pt}\end{center}

\section{Chủ đề bổ sung 3.7 -- Chương trình đồ họa tương
tác}\label{chux1ee7-ux111ux1ec1-bux1ed5-sung-3.7-chux1b0ux1a1ng-truxecnh-ux111ux1ed3-hux1ecda-tux1b0ux1a1ng-tuxe1c}

Logic rẽ nhánh cũng áp dụng trong chương trình đồ họa.

Một chương trình có thể:

\begin{enumerate}
\def\labelenumi{\arabic{enumi}.}
\tightlist
\item
  đọc dữ liệu;
\item
  kiểm tra tọa độ hoặc kích thước;
\item
  chỉ vẽ khi dữ liệu hợp lệ;
\item
  phản ứng với một cú nhấp chuột và rẽ nhánh theo vị trí nhấp.
\end{enumerate}

Điểm quan trọng là giao diện đồ họa không làm mất nhu cầu kiểm tra dữ
liệu đầu vào.

\begin{center}\rule{0.5\linewidth}{0.5pt}\end{center}

\section{Máy tính và xã hội 3.2 -- Trí tuệ nhân
tạo}\label{muxe1y-tuxednh-vuxe0-xuxe3-hux1ed9i-3.2-truxed-tuux1ec7-nhuxe2n-tux1ea1o}

Chương giới thiệu trí tuệ nhân tạo như một bối cảnh rộng hơn cho phần
mềm ra quyết định. Chương trình truyền thống tuân theo các quy tắc do
lập trình viên mô tả rõ, trong khi AI hướng đến những hệ thống có hành
vi linh hoạt hoặc thông minh hơn.

Liên hệ với bài học này chủ yếu mang tính khái niệm: ngay cả các hệ
thống phức tạp vẫn phải biểu diễn thông tin, đánh giá phương án và chọn
hành động. Các kỹ thuật AI hiện đại có thể học một phần quy tắc quyết
định từ dữ liệu, nhưng hiểu luồng điều khiển tường minh vẫn là nền tảng
của lập trình.

\begin{center}\rule{0.5\linewidth}{0.5pt}\end{center}

\section{Ví dụ minh họa 3.2 -- Hai đường tròn giao
nhau}\label{vuxed-dux1ee5-minh-hux1ecda-3.2-hai-ux111ux1b0ux1eddng-truxf2n-giao-nhau}

Ví dụ đồ họa của chương xác định quan hệ giữa hai đường tròn.

Giả sử tâm hai đường tròn là:

\begin{Shaded}
\begin{Highlighting}[]
\NormalTok{(x0, y0)}
\NormalTok{(x1, y1)}
\end{Highlighting}
\end{Shaded}

và bán kính:

\begin{Shaded}
\begin{Highlighting}[]
\NormalTok{r0}
\NormalTok{r1}
\end{Highlighting}
\end{Shaded}

Khoảng cách giữa hai tâm:

\begin{Shaded}
\begin{Highlighting}[]
\ImportTok{from}\NormalTok{ math }\ImportTok{import}\NormalTok{ sqrt}

\NormalTok{dist }\OperatorTok{=}\NormalTok{ sqrt((x1 }\OperatorTok{{-}}\NormalTok{ x0) }\OperatorTok{**} \DecValTok{2} \OperatorTok{+}\NormalTok{ (y1 }\OperatorTok{{-}}\NormalTok{ y0) }\OperatorTok{**} \DecValTok{2}\NormalTok{)}
\end{Highlighting}
\end{Shaded}

Có thể phân loại quan hệ bằng rẽ nhánh:

\begin{Shaded}
\begin{Highlighting}[]
\ControlFlowTok{if}\NormalTok{ dist }\OperatorTok{\textgreater{}}\NormalTok{ r0 }\OperatorTok{+}\NormalTok{ r1 :}
\NormalTok{    message }\OperatorTok{=} \StringTok{"The circles are separate."}
\ControlFlowTok{elif}\NormalTok{ dist }\OperatorTok{\textless{}} \BuiltInTok{abs}\NormalTok{(r0 }\OperatorTok{{-}}\NormalTok{ r1) :}
\NormalTok{    message }\OperatorTok{=} \StringTok{"One circle is inside the other."}
\ControlFlowTok{elif}\NormalTok{ dist }\OperatorTok{==}\NormalTok{ r0 }\OperatorTok{+}\NormalTok{ r1 :}
\NormalTok{    message }\OperatorTok{=} \StringTok{"The circles touch externally."}
\ControlFlowTok{elif}\NormalTok{ dist }\OperatorTok{==} \DecValTok{0} \KeywordTok{and}\NormalTok{ r0 }\OperatorTok{==}\NormalTok{ r1 :}
\NormalTok{    message }\OperatorTok{=} \StringTok{"The circles coincide."}
\ControlFlowTok{else}\NormalTok{ :}
\NormalTok{    message }\OperatorTok{=} \StringTok{"The circles intersect at two points."}
\end{Highlighting}
\end{Shaded}

Ví dụ này kết hợp:

\begin{itemize}
\tightlist
\item
  kiểm tra đầu vào;
\item
  số học;
\item
  toán tử quan hệ;
\item
  toán tử Boolean;
\item
  nhiều phương án;
\item
  hình học;
\item
  đầu ra đồ họa.
\end{itemize}

Với các phép tính hình học dùng \texttt{float}, nên cân nhắc kiểm tra
gần bằng thay vì bằng tuyệt đối.

\begin{center}\rule{0.5\linewidth}{0.5pt}\end{center}

\section{TOOLBOX 3.2 -- Vẽ đồ thị đơn
giản}\label{toolbox-3.2-vux1ebd-ux111ux1ed3-thux1ecb-ux111ux1a1n-giux1ea3n}

Chương giới thiệu \texttt{matplotlib.pyplot} như một ví dụ tùy chọn về
thư viện Python dùng để trực quan hóa dữ liệu.

Biểu đồ cột đơn giản:

\begin{Shaded}
\begin{Highlighting}[]
\ImportTok{from}\NormalTok{ matplotlib }\ImportTok{import}\NormalTok{ pyplot}

\NormalTok{pyplot.bar(}\DecValTok{1}\NormalTok{, }\FloatTok{1.1}\NormalTok{)}
\NormalTok{pyplot.bar(}\DecValTok{2}\NormalTok{, }\FloatTok{10.0}\NormalTok{)}
\NormalTok{pyplot.bar(}\DecValTok{3}\NormalTok{, }\FloatTok{25.4}\NormalTok{)}
\NormalTok{pyplot.bar(}\DecValTok{4}\NormalTok{, }\FloatTok{44.5}\NormalTok{)}
\NormalTok{pyplot.bar(}\DecValTok{5}\NormalTok{, }\FloatTok{61.0}\NormalTok{)}

\NormalTok{pyplot.xlabel(}\StringTok{"Month"}\NormalTok{)}
\NormalTok{pyplot.ylabel(}\StringTok{"Temperature"}\NormalTok{)}
\NormalTok{pyplot.show()}
\end{Highlighting}
\end{Shaded}

Đồ thị đường đơn giản:

\begin{Shaded}
\begin{Highlighting}[]
\ImportTok{from}\NormalTok{ matplotlib }\ImportTok{import}\NormalTok{ pyplot}

\NormalTok{months }\OperatorTok{=}\NormalTok{ [}\DecValTok{1}\NormalTok{, }\DecValTok{2}\NormalTok{, }\DecValTok{3}\NormalTok{, }\DecValTok{4}\NormalTok{, }\DecValTok{5}\NormalTok{]}
\NormalTok{temperatures }\OperatorTok{=}\NormalTok{ [}\FloatTok{1.1}\NormalTok{, }\FloatTok{10.0}\NormalTok{, }\FloatTok{25.4}\NormalTok{, }\FloatTok{44.5}\NormalTok{, }\FloatTok{61.0}\NormalTok{]}

\NormalTok{pyplot.plot(months, temperatures)}
\NormalTok{pyplot.xlabel(}\StringTok{"Month"}\NormalTok{)}
\NormalTok{pyplot.ylabel(}\StringTok{"Temperature"}\NormalTok{)}
\NormalTok{pyplot.show()}
\end{Highlighting}
\end{Shaded}

Toolbox về đồ thị là phần tùy chọn. Mục tiêu chính của chương vẫn là câu
lệnh rẽ nhánh và logic Boolean.

\begin{center}\rule{0.5\linewidth}{0.5pt}\end{center}

\section{Tổng kết Bài 3}\label{tux1ed5ng-kux1ebft-buxe0i-3}

\subsection{\texorpdfstring{Câu lệnh
\texttt{if}}{Câu lệnh if}}\label{cuxe2u-lux1ec7nh-if-1}

\begin{itemize}
\tightlist
\item
  \texttt{if} thay đổi đường thực hiện của chương trình dựa trên một
  điều kiện.
\item
  \texttt{else} cung cấp phương án thay thế khi điều kiện sai.
\item
  \texttt{elif} hỗ trợ nhiều phương án loại trừ lẫn nhau.
\item
  Câu lệnh ghép có dòng tiêu đề và một khối thụt lề.
\item
  Thụt lề nhất quán là một phần của cú pháp Python.
\item
  Tránh lặp mã giữa các nhánh khi phần chung có thể đưa ra ngoài.
\end{itemize}

\subsection{Các toán tử quan
hệ}\label{cuxe1c-touxe1n-tux1eed-quan-hux1ec7-2}

\begin{Shaded}
\begin{Highlighting}[]
\OperatorTok{\textless{}}  \OperatorTok{\textless{}=}  \OperatorTok{\textgreater{}}  \OperatorTok{\textgreater{}=}  \OperatorTok{==}  \OperatorTok{!=}
\end{Highlighting}
\end{Shaded}

Ghi nhớ:

\begin{Shaded}
\begin{Highlighting}[]
\NormalTok{=   gán}
\NormalTok{==  kiểm tra bằng nhau}
\end{Highlighting}
\end{Shaded}

Chuỗi cũng có thể được so sánh và phép so sánh mặc định phân biệt chữ
hoa/chữ thường.

Số thực được tạo từ phép tính thường nên so sánh bằng sai số cho phép
thay vì bằng tuyệt đối.

\subsection{Rẽ nhánh lồng nhau}\label{rux1ebd-nhuxe1nh-lux1ed3ng-nhau-1}

\begin{itemize}
\tightlist
\item
  Một câu lệnh \texttt{if} có thể nằm trong một nhánh khác.
\item
  Rẽ nhánh lồng nhau phù hợp với quyết định nhiều tầng.
\item
  Thụt lề cho biết quyết định bên trong thuộc nhánh nào.
\item
  Mô phỏng bằng tay giúp kiểm tra logic lồng nhau.
\end{itemize}

\subsection{Nhiều phương án}\label{nhiux1ec1u-phux1b0ux1a1ng-uxe1n-2}

Dùng:

\begin{Shaded}
\begin{Highlighting}[]
\ControlFlowTok{if}\NormalTok{ ... :}
\NormalTok{    ...}
\ControlFlowTok{elif}\NormalTok{ ... :}
\NormalTok{    ...}
\ControlFlowTok{else}\NormalTok{ :}
\NormalTok{    ...}
\end{Highlighting}
\end{Shaded}

Chỉ nhánh đúng đầu tiên được thực hiện.

Với các ngưỡng chồng lấn, phải sắp xếp điều kiện cẩn thận.

\subsection{Lưu đồ}\label{lux1b0u-ux111ux1ed3}

Lưu đồ biểu diễn trực quan:

\begin{itemize}
\tightlist
\item
  công việc xử lý;
\item
  nhập/xuất;
\item
  quyết định;
\item
  nhánh và điểm nhập lại.
\end{itemize}

Nên dùng cấu trúc rõ ràng và tránh các mũi tên ``nhảy'' vào giữa nhánh
khác.

\subsection{Ca kiểm thử}\label{ca-kiux1ec3m-thux1eed}

Ca kiểm thử tốt nên gồm:

\begin{itemize}
\tightlist
\item
  ít nhất một ca cho mỗi nhánh;
\item
  giá trị biên;
\item
  các giá trị ngay bên dưới và bên trên biên;
\item
  giá trị không hợp lệ khi có yêu cầu validation.
\end{itemize}

Thiết kế kết quả mong đợi trước khi viết mã giúp cả việc xây dựng thuật
toán lẫn kiểm thử.

\subsection{Logic Boolean}\label{logic-boolean}

\begin{Shaded}
\begin{Highlighting}[]
\VariableTok{True}
\VariableTok{False}
\KeywordTok{and}
\KeywordTok{or}
\KeywordTok{not}
\end{Highlighting}
\end{Shaded}

\begin{itemize}
\tightlist
\item
  \texttt{and} yêu cầu tất cả điều kiện thành phần đúng.
\item
  \texttt{or} yêu cầu ít nhất một điều kiện đúng.
\item
  \texttt{not} đảo giá trị Boolean.
\item
  Python sử dụng đánh giá ngắn mạch.
\item
  Định luật De Morgan giúp biến đổi điều kiện phủ định phức tạp.
\end{itemize}

\subsection{Phân tích chuỗi}\label{phuxe2n-tuxedch-chuux1ed7i-1}

Các phép toán hữu ích:

\begin{Shaded}
\begin{Highlighting}[]
\NormalTok{substring }\KeywordTok{in}\NormalTok{ text}
\NormalTok{substring }\KeywordTok{not} \KeywordTok{in}\NormalTok{ text}
\NormalTok{text.count(...)}
\NormalTok{text.find(...)}
\NormalTok{text.startswith(...)}
\NormalTok{text.endswith(...)}
\NormalTok{text.isalpha()}
\NormalTok{text.isdigit()}
\NormalTok{text.isalnum()}
\NormalTok{text.islower()}
\NormalTok{text.isupper()}
\NormalTok{text.isspace()}
\end{Highlighting}
\end{Shaded}

\subsection{Kiểm tra dữ liệu đầu
vào}\label{kiux1ec3m-tra-dux1eef-liux1ec7u-ux111ux1ea7u-vuxe0o}

Trước khi xử lý dữ liệu do người dùng nhập:

\begin{Shaded}
\begin{Highlighting}[]
\NormalTok{Đọc}
\NormalTok{→ chuẩn hóa nếu cần}
\NormalTok{→ kiểm tra hợp lệ}
\NormalTok{→ từ chối giá trị sai}
\NormalTok{→ xử lý giá trị đúng}
\end{Highlighting}
\end{Shaded}

Kiểm tra miền giá trị và các lựa chọn hợp lệ có thể làm bằng điều kiện.
Lỗi chuyển đổi số cần xử lý ngoại lệ, sẽ học sau.

\begin{center}\rule{0.5\linewidth}{0.5pt}\end{center}

\section{Quiz tổng kết}\label{quiz-tux1ed5ng-kux1ebft}

\subsection{Câu 1}\label{cuxe2u-1-8}

Toán tử nào kiểm tra bằng nhau?

A. \texttt{=}\\
B. \texttt{==}\\
C. \texttt{!=}\\
D. \texttt{\textless{}=}

\textbf{Đáp án}

\textbf{B. \texttt{==}}

\subsection{Câu 2}\label{cuxe2u-2-8}

Điều gì xảy ra khi điều kiện của \texttt{if/else} là đúng?

A. Cả hai nhánh đều chạy.\\
B. Chỉ nhánh \texttt{else} chạy.\\
C. Nhánh \texttt{if} chạy và \texttt{else} bị bỏ qua.\\
D. Chương trình luôn kết thúc.

\textbf{Đáp án}

\textbf{C}

\subsection{Câu 3}\label{cuxe2u-3-5}

Giá trị của:

\begin{Shaded}
\begin{Highlighting}[]
\DecValTok{5} \OperatorTok{\textgreater{}=} \DecValTok{5}
\end{Highlighting}
\end{Shaded}

là gì?

A. \texttt{True}\\
B. \texttt{False}\\
C. \texttt{5}\\
D. Lỗi

\textbf{Đáp án}

\textbf{A. \texttt{True}}

\subsection{Câu 4}\label{cuxe2u-4-3}

Điều kiện nào kiểm tra \texttt{x} nằm \emph{nghiêm ngặt} giữa \texttt{0}
và \texttt{10}?

A. \texttt{x\ \textgreater{}\ 0\ or\ x\ \textless{}\ 10}\\
B. \texttt{x\ \textgreater{}\ 0\ and\ x\ \textless{}\ 10}\\
C. \texttt{x\ \textgreater{}=\ 0\ or\ x\ \textless{}=\ 10}\\
D. \texttt{not\ x}

\textbf{Đáp án}

\textbf{B}

\subsection{Câu 5}\label{cuxe2u-5}

Kết quả của:

\begin{Shaded}
\begin{Highlighting}[]
\VariableTok{True} \KeywordTok{or} \VariableTok{False}
\end{Highlighting}
\end{Shaded}

là gì?

\textbf{Đáp án}

\texttt{True}.

\subsection{Câu 6}\label{cuxe2u-6}

Tại sao phép kiểm tra sau có thể không an toàn sau các phép tính số
thực?

\begin{Shaded}
\begin{Highlighting}[]
\ControlFlowTok{if}\NormalTok{ result }\OperatorTok{==}\NormalTok{ expected :}
\end{Highlighting}
\end{Shaded}

\textbf{Đáp án}

Vì sai số làm tròn nhỏ có thể khiến hai giá trị về mặt toán học bằng
nhau nhưng biểu diễn \texttt{float} hơi khác nhau. Trong nhiều trường
hợp nên so sánh theo một ngưỡng sai số.

\subsection{Câu 7}\label{cuxe2u-7}

Cấu trúc nào thường phù hợp nhất cho các loại điểm loại trừ lẫn nhau?

A. Nhiều \texttt{if} độc lập\\
B. \texttt{if/elif/else}\\
C. Chỉ \texttt{else}\\
D. Không cần điều kiện

\textbf{Đáp án}

\textbf{B. \texttt{if/elif/else}}

\subsection{Câu 8}\label{cuxe2u-8}

Biểu thức sau kiểm tra điều gì?

\begin{Shaded}
\begin{Highlighting}[]
\CommentTok{".csv"} \KeywordTok{in}\NormalTok{ filename}
\end{Highlighting}
\end{Shaded}

\textbf{Đáp án}

Kiểm tra chuỗi con chính xác \texttt{".csv"} có xuất hiện ở bất kỳ vị
trí nào trong \texttt{filename} hay không.

\subsection{Câu 9}\label{cuxe2u-9}

\texttt{text.find("abc")} trả về gì khi không có \texttt{"abc"}?

A. \texttt{False}\\
B. \texttt{None}\\
C. \texttt{-1}\\
D. Lỗi

\textbf{Đáp án}

\textbf{C. \texttt{-1}}

\subsection{Câu 10}\label{cuxe2u-10}

Giá trị biên cho:

\begin{Shaded}
\begin{Highlighting}[]
\ControlFlowTok{if}\NormalTok{ score }\OperatorTok{\textgreater{}=} \DecValTok{50}\NormalTok{ :}
\end{Highlighting}
\end{Shaded}

là gì?

\textbf{Đáp án}

\texttt{50} là giá trị biên quan trọng. \texttt{49} và \texttt{51} là
hai giá trị lân cận hữu ích.

\subsection{Câu 11}\label{cuxe2u-11}

Tại sao thứ tự điều kiện trong chuỗi \texttt{if/elif} có thể ảnh hưởng
kết quả?

\textbf{Đáp án}

Vì Python dừng ở nhánh đúng đầu tiên. Một điều kiện quá rộng đặt quá sớm
có thể khiến điều kiện cụ thể hơn ở phía sau không bao giờ được kiểm
tra.

\subsection{Câu 12}\label{cuxe2u-12}

Mục tiêu chính của kiểm tra dữ liệu đầu vào là gì?

A. Làm mã dài hơn.\\
B. Đảm bảo dữ liệu người dùng nhập thỏa các giả định của chương trình
trước khi xử lý.\\
C. Chuyển mọi giá trị thành chuỗi.\\
D. Loại bỏ tất cả câu lệnh \texttt{if}.

\textbf{Đáp án}

\textbf{B}

\begin{center}\rule{0.5\linewidth}{0.5pt}\end{center}

\section{Bài tập ôn tập}\label{buxe0i-tux1eadp-uxf4n-tux1eadp}

\subsection{Bài 1. Theo dõi quyết định hai
nhánh}\label{buxe0i-1.-theo-duxf5i-quyux1ebft-ux111ux1ecbnh-hai-nhuxe1nh}

Xác định giá trị cuối cùng của \texttt{fee}:

\begin{Shaded}
\begin{Highlighting}[]
\NormalTok{age }\OperatorTok{=} \DecValTok{16}

\ControlFlowTok{if}\NormalTok{ age }\OperatorTok{\textless{}} \DecValTok{18}\NormalTok{ :}
\NormalTok{    fee }\OperatorTok{=} \DecValTok{5}
\ControlFlowTok{else}\NormalTok{ :}
\NormalTok{    fee }\OperatorTok{=} \DecValTok{10}
\end{Highlighting}
\end{Shaded}

\textbf{Đáp án}

\texttt{fee\ =\ 5}.

\begin{center}\rule{0.5\linewidth}{0.5pt}\end{center}

\subsection{Bài 2. Chọn đúng toán tử quan
hệ}\label{buxe0i-2.-chux1ecdn-ux111uxfang-touxe1n-tux1eed-quan-hux1ec7}

Điền toán tử để thông báo được in khi \texttt{temperature} không lớn hơn
\texttt{0}.

\begin{Shaded}
\begin{Highlighting}[]
\ControlFlowTok{if}\NormalTok{ temperature \_\_\_ }\DecValTok{0}\NormalTok{ :}
    \BuiltInTok{print}\NormalTok{(}\StringTok{"Freezing or below"}\NormalTok{)}
\end{Highlighting}
\end{Shaded}

\textbf{Đáp án}

\begin{Shaded}
\begin{Highlighting}[]
\OperatorTok{\textless{}=}
\end{Highlighting}
\end{Shaded}

\begin{center}\rule{0.5\linewidth}{0.5pt}\end{center}

\subsection{Bài 3. Viết lại để tránh lặp
mã}\label{buxe0i-3.-viux1ebft-lux1ea1i-ux111ux1ec3-truxe1nh-lux1eb7p-muxe3}

Cải thiện đoạn sau:

\begin{Shaded}
\begin{Highlighting}[]
\ControlFlowTok{if}\NormalTok{ quantity }\OperatorTok{\textgreater{}=} \DecValTok{10}\NormalTok{ :}
\NormalTok{    rate }\OperatorTok{=} \FloatTok{0.90}
\NormalTok{    total }\OperatorTok{=}\NormalTok{ rate }\OperatorTok{*}\NormalTok{ price }\OperatorTok{*}\NormalTok{ quantity}
\ControlFlowTok{else}\NormalTok{ :}
\NormalTok{    rate }\OperatorTok{=} \FloatTok{1.00}
\NormalTok{    total }\OperatorTok{=}\NormalTok{ rate }\OperatorTok{*}\NormalTok{ price }\OperatorTok{*}\NormalTok{ quantity}
\end{Highlighting}
\end{Shaded}

\textbf{Đáp án tham khảo}

\begin{Shaded}
\begin{Highlighting}[]
\ControlFlowTok{if}\NormalTok{ quantity }\OperatorTok{\textgreater{}=} \DecValTok{10}\NormalTok{ :}
\NormalTok{    rate }\OperatorTok{=} \FloatTok{0.90}
\ControlFlowTok{else}\NormalTok{ :}
\NormalTok{    rate }\OperatorTok{=} \FloatTok{1.00}

\NormalTok{total }\OperatorTok{=}\NormalTok{ rate }\OperatorTok{*}\NormalTok{ price }\OperatorTok{*}\NormalTok{ quantity}
\end{Highlighting}
\end{Shaded}

\begin{center}\rule{0.5\linewidth}{0.5pt}\end{center}

\subsection{Bài 4. Thiết kế ca kiểm
thử}\label{buxe0i-4.-thiux1ebft-kux1ebf-ca-kiux1ec3m-thux1eed}

Với:

\begin{Shaded}
\begin{Highlighting}[]
\ControlFlowTok{if}\NormalTok{ balance }\OperatorTok{\textless{}} \DecValTok{0}\NormalTok{ :}
\NormalTok{    status }\OperatorTok{=} \StringTok{"Overdrawn"}
\ControlFlowTok{else}\NormalTok{ :}
\NormalTok{    status }\OperatorTok{=} \StringTok{"OK"}
\end{Highlighting}
\end{Shaded}

hãy đề xuất ba giá trị kiểm thử hữu ích.

\textbf{Đáp án tham khảo}

\begin{Shaded}
\begin{Highlighting}[]
\NormalTok{{-}1   → Overdrawn}
\NormalTok{0    → OK       (biên)}
\NormalTok{1    → OK}
\end{Highlighting}
\end{Shaded}

\begin{center}\rule{0.5\linewidth}{0.5pt}\end{center}

\subsection{Bài 5. Biểu thức
Boolean}\label{buxe0i-5.-biux1ec3u-thux1ee9c-boolean}

Viết điều kiện cho các trường hợp:

\begin{enumerate}
\def\labelenumi{\arabic{enumi}.}
\tightlist
\item
  \texttt{x} và \texttt{y} đều dương.
\item
  Ít nhất một trong \texttt{x}, \texttt{y} âm.
\item
  \texttt{age} nằm trong khoảng từ \texttt{18} đến \texttt{65}, kể cả
  hai đầu.
\item
  \texttt{answer} không phải \texttt{"y"} cũng không phải \texttt{"n"}.
\end{enumerate}

\textbf{Đáp án tham khảo}

\begin{Shaded}
\begin{Highlighting}[]
\NormalTok{x }\OperatorTok{\textgreater{}} \DecValTok{0} \KeywordTok{and}\NormalTok{ y }\OperatorTok{\textgreater{}} \DecValTok{0}
\NormalTok{x }\OperatorTok{\textless{}} \DecValTok{0} \KeywordTok{or}\NormalTok{ y }\OperatorTok{\textless{}} \DecValTok{0}
\NormalTok{age }\OperatorTok{\textgreater{}=} \DecValTok{18} \KeywordTok{and}\NormalTok{ age }\OperatorTok{\textless{}=} \DecValTok{65}
\NormalTok{answer }\OperatorTok{!=} \StringTok{"y"} \KeywordTok{and}\NormalTok{ answer }\OperatorTok{!=} \StringTok{"n"}
\end{Highlighting}
\end{Shaded}

\begin{center}\rule{0.5\linewidth}{0.5pt}\end{center}

\subsection{Bài 6. Phân tích
chuỗi}\label{buxe0i-6.-phuxe2n-tuxedch-chuux1ed7i}

Cho:

\begin{Shaded}
\begin{Highlighting}[]
\NormalTok{filename }\OperatorTok{=} \StringTok{"report.Final.PDF"}
\end{Highlighting}
\end{Shaded}

viết điều kiện không phân biệt chữ hoa/thường để kiểm tra file có kết
thúc bằng \texttt{.pdf} hay không.

\textbf{Đáp án}

\begin{Shaded}
\begin{Highlighting}[]
\NormalTok{filename.lower().endswith(}\StringTok{".pdf"}\NormalTok{)}
\end{Highlighting}
\end{Shaded}

\begin{center}\rule{0.5\linewidth}{0.5pt}\end{center}

\section{Bài tập thực hành}\label{buxe0i-tux1eadp-thux1ef1c-huxe0nh}

\subsection{Bài 1. Đạt / Không đạt / Điểm không hợp
lệ}\label{buxe0i-1.-ux111ux1ea1t-khuxf4ng-ux111ux1ea1t-ux111iux1ec3m-khuxf4ng-hux1ee3p-lux1ec7}

Nhập một điểm thi.

\begin{itemize}
\tightlist
\item
  Nếu điểm nhỏ hơn \texttt{0} hoặc lớn hơn \texttt{100}, in thông báo
  lỗi.
\item
  Nếu hợp lệ, in \texttt{"Pass"} cho điểm từ \texttt{50} trở lên và
  \texttt{"Fail"} cho điểm thấp hơn.
\end{itemize}

Hãy vẽ lưu đồ trước khi viết Python.

\begin{center}\rule{0.5\linewidth}{0.5pt}\end{center}

\subsection{Bài 2. Số lớn hơn trong hai
số}\label{buxe0i-2.-sux1ed1-lux1edbn-hux1a1n-trong-hai-sux1ed1}

Nhập hai số và in:

\begin{itemize}
\tightlist
\item
  số lớn hơn;
\item
  \texttt{"Equal"} nếu hai số bằng nhau.
\end{itemize}

Dùng \texttt{if/elif/else}.

\begin{center}\rule{0.5\linewidth}{0.5pt}\end{center}

\subsection{Bài 3. Số lớn nhất trong ba
số}\label{buxe0i-3.-sux1ed1-lux1edbn-nhux1ea5t-trong-ba-sux1ed1}

Nhập ba số và xác định số lớn nhất.

Trước tiên giải bằng rẽ nhánh lồng nhau. Sau đó xem có thể viết rõ hơn
bằng biểu thức Boolean hoặc nhiều phương án hay không.

\begin{center}\rule{0.5\linewidth}{0.5pt}\end{center}

\subsection{Bài 4. Quy tắc vận
chuyển}\label{buxe0i-4.-quy-tux1eafc-vux1eadn-chuyux1ec3n}

Nhập:

\begin{itemize}
\tightlist
\item
  quốc gia đích;
\item
  bang hoặc tỉnh.
\end{itemize}

Chuẩn hóa hai đầu vào bằng \texttt{upper()}.

Dùng quy tắc:

\begin{itemize}
\tightlist
\item
  phần lục địa Hoa Kỳ: \texttt{\$5};
\item
  Alaska hoặc Hawaii: \texttt{\$10};
\item
  ngoài Hoa Kỳ: \texttt{\$15}.
\end{itemize}

In phí vận chuyển với hai chữ số thập phân.

\begin{center}\rule{0.5\linewidth}{0.5pt}\end{center}

\subsection{Bài 5. Phân loại điểm
chữ}\label{buxe0i-5.-phuxe2n-loux1ea1i-ux111iux1ec3m-chux1eef}

Nhập điểm từ \texttt{0} đến \texttt{100} và phân loại:

\begin{Shaded}
\begin{Highlighting}[]
\NormalTok{90–100 → A}
\NormalTok{80–89  → B}
\NormalTok{70–79  → C}
\NormalTok{60–69  → D}
\NormalTok{0–59   → F}
\end{Highlighting}
\end{Shaded}

Từ chối điểm không hợp lệ trước khi phân loại.

Tạo ca kiểm thử cho:

\begin{Shaded}
\begin{Highlighting}[]
\NormalTok{{-}1, 0, 59, 60, 69, 70, 79, 80, 89, 90, 100, 101}
\end{Highlighting}
\end{Shaded}

\begin{center}\rule{0.5\linewidth}{0.5pt}\end{center}

\subsection{Bài 6. Kiểm tra tên người
dùng}\label{buxe0i-6.-kiux1ec3m-tra-tuxean-ngux1b0ux1eddi-duxf9ng}

Nhập một username và kiểm tra:

\begin{itemize}
\tightlist
\item
  không được rỗng;
\item
  chỉ gồm chữ cái và chữ số;
\item
  có ít nhất một ký tự.
\end{itemize}

Dùng các phương thức phân tích chuỗi ở Mục 3.8.

\begin{center}\rule{0.5\linewidth}{0.5pt}\end{center}

\subsection{Bài 7. Phân loại loại
file}\label{buxe0i-7.-phuxe2n-loux1ea1i-loux1ea1i-file}

Nhập tên file và phân loại:

\begin{itemize}
\tightlist
\item
  mã Python: \texttt{.py};
\item
  notebook: \texttt{.ipynb};
\item
  dữ liệu CSV: \texttt{.csv};
\item
  file văn bản: \texttt{.txt};
\item
  loại khác.
\end{itemize}

Kiểm tra không phân biệt chữ hoa/chữ thường.

\begin{center}\rule{0.5\linewidth}{0.5pt}\end{center}

\subsection{Bài 8. So sánh gần
bằng}\label{buxe0i-8.-so-suxe1nh-gux1ea7n-bux1eb1ng}

Nhập hai số thực và xác định chúng có chênh nhau ít hơn \texttt{0.001}
hay không.

Không dùng phép bằng tuyệt đối cho phép so sánh gần bằng.

\begin{center}\rule{0.5\linewidth}{0.5pt}\end{center}

\subsection{Bài 9. Phân loại tam giác theo độ dài
cạnh}\label{buxe0i-9.-phuxe2n-loux1ea1i-tam-giuxe1c-theo-ux111ux1ed9-duxe0i-cux1ea1nh}

Nhập ba độ dài cạnh dương.

Trước tiên kiểm tra chúng có tạo thành tam giác hay không:

\begin{Shaded}
\begin{Highlighting}[]
\NormalTok{a + b \textgreater{} c}
\NormalTok{and}
\NormalTok{b + c \textgreater{} a}
\NormalTok{and}
\NormalTok{c + a \textgreater{} b}
\end{Highlighting}
\end{Shaded}

Nếu hợp lệ, phân loại:

\begin{itemize}
\tightlist
\item
  tam giác đều;
\item
  tam giác cân;
\item
  tam giác thường.
\end{itemize}

Chuẩn bị ca kiểm thử cho từng nhánh và các giá trị biên trước khi viết
mã.

\begin{center}\rule{0.5\linewidth}{0.5pt}\end{center}

\subsection{Bài 10. Kiểm tra mã truy cập đơn
giản}\label{buxe0i-10.-kiux1ec3m-tra-muxe3-truy-cux1eadp-ux111ux1a1n-giux1ea3n}

Nhập một mã truy cập ngắn.

Quy tắc:

\begin{itemize}
\tightlist
\item
  đúng 6 ký tự;
\item
  chỉ gồm chữ và số;
\item
  phải chứa chuỗi con \texttt{"NEU"}, không phân biệt hoa/thường.
\end{itemize}

In thông báo cụ thể cho trường hợp hợp lệ hoặc không hợp lệ.

\begin{center}\rule{0.5\linewidth}{0.5pt}\end{center}

\section{Bài tập mở}\label{buxe0i-tux1eadp-mux1edf}

\subsection{Bài 1}\label{buxe0i-1}

Đưa ra ba ví dụ thực tế trong đó một quyết định chỉ có đúng hai kết quả.
Viết quy tắc cho mỗi ví dụ bằng mã giả.

\subsection{Bài 2}\label{buxe0i-2}

Đưa ra một ví dụ thực tế có ít nhất bốn phương án loại trừ lẫn nhau.
Giải thích vì sao \texttt{if/elif/else} phù hợp.

\subsection{Bài 3}\label{buxe0i-3}

Chọn một điều kiện dạng khoảng và xác định:

\begin{itemize}
\tightlist
\item
  một giá trị hợp lệ điển hình;
\item
  giá trị biên dưới;
\item
  giá trị biên trên;
\item
  một giá trị ngay dưới khoảng;
\item
  một giá trị ngay trên khoảng.
\end{itemize}

\subsection{Bài 4}\label{buxe0i-4}

Tạo ví dụ trong đó nhiều \texttt{if} độc lập là đúng nhưng
\texttt{if/elif} lại sai. Giải thích vì sao có thể cần thực hiện nhiều
hành động cùng lúc.

\subsection{Bài 5}\label{buxe0i-5}

Viết một điều kiện Boolean phức tạp dùng \texttt{not}, \texttt{and},
\texttt{or}. Sau đó đơn giản hóa bằng định luật De Morgan.

\begin{center}\rule{0.5\linewidth}{0.5pt}\end{center}

\section{Thuật ngữ chính}\label{thuux1eadt-ngux1eef-chuxednh}

\begin{longtable}[]{@{}
  >{\raggedright\arraybackslash}p{(\columnwidth - 2\tabcolsep) * \real{0.5000}}
  >{\raggedright\arraybackslash}p{(\columnwidth - 2\tabcolsep) * \real{0.5000}}@{}}
\toprule\noalign{}
\begin{minipage}[b]{\linewidth}\raggedright
Thuật ngữ
\end{minipage} & \begin{minipage}[b]{\linewidth}\raggedright
Ý nghĩa
\end{minipage} \\
\midrule\noalign{}
\endhead
\bottomrule\noalign{}
\endlastfoot
Decision & Quyết định -- chọn hành động dựa trên điều kiện \\
Condition & Điều kiện -- biểu thức cho kết quả đúng hoặc sai \\
\texttt{if} statement & Câu lệnh \texttt{if} -- thực hiện một khối khi
điều kiện đúng \\
\texttt{else} branch & Nhánh thay thế khi điều kiện \texttt{if} sai \\
\texttt{elif} & Thêm một điều kiện khác vào quyết định nhiều nhánh \\
Compound statement & Câu lệnh ghép gồm dòng tiêu đề và các câu lệnh lồng
bên dưới \\
Statement block & Nhóm câu lệnh ở cùng mức thụt lề \\
Indentation & Khoảng trắng đầu dòng xác định cấu trúc khối trong
Python \\
Relational operator & Toán tử quan hệ so sánh hai giá trị \\
Equality test & Kiểm tra bằng nhau bằng \texttt{==} \\
Boundary value & Giá trị nằm đúng tại biên của một khoảng quyết định \\
Floating-point tolerance & Sai số cho phép dùng khi so sánh gần bằng \\
Nested branch & Câu lệnh quyết định nằm trong một nhánh khác \\
Hand-tracing & Mô phỏng thủ công quá trình chạy chương trình \\
Multiple alternatives & Quyết định có nhiều hơn hai nhánh \\
Flowchart & Sơ đồ mô tả công việc và luồng điều khiển \\
Branch coverage & Kiểm thử để mọi nhánh đều được chạy ít nhất một lần \\
Test case & Đầu vào cùng kết quả mong đợi để kiểm chứng chương trình \\
Boolean & Kiểu dữ liệu có hai giá trị \texttt{True} và \texttt{False} \\
Flag & Biến Boolean biểu diễn trạng thái hoặc điều kiện \\
Boolean operator & \texttt{and}, \texttt{or}, \texttt{not} \\
Short-circuit evaluation & Dừng đánh giá Boolean khi đã biết kết quả
cuối \\
De Morgan's laws & Quy tắc biến đổi phủ định của biểu thức
\texttt{and}/\texttt{or} \\
Substring & Chuỗi con nằm trong một chuỗi khác \\
Membership test & Kiểm tra bằng \texttt{in} hoặc \texttt{not\ in} \\
Input validation & Kiểm tra dữ liệu đầu vào thỏa các ràng buộc \\
Normalization & Chuyển dữ liệu về một dạng thống nhất trước khi so
sánh \\
\end{longtable}

\begin{center}\rule{0.5\linewidth}{0.5pt}\end{center}

\section{Hướng dẫn học
tập}\label{hux1b0ux1edbng-dux1eabn-hux1ecdc-tux1eadp}

Sau Bài 3, sinh viên không nên còn xem chương trình chỉ là một danh sách
phép tính cố định. Chương trình giờ đã có thể \textbf{chọn đường thực
hiện}.

Một quy trình phù hợp cho bài toán có rẽ nhánh là:

\begin{Shaded}
\begin{Highlighting}[]
\NormalTok{Hiểu quy tắc}
\NormalTok{→ xác định các kết quả có thể xảy ra}
\NormalTok{→ xây dựng điều kiện}
\NormalTok{→ kiểm tra giá trị biên}
\NormalTok{→ vẽ lưu đồ khi cần}
\NormalTok{→ viết mã giả}
\NormalTok{→ chọn if / if{-}else / if{-}elif{-}else / if lồng nhau}
\NormalTok{→ loại bỏ phần xử lý lặp lại}
\NormalTok{→ chuẩn bị ca kiểm thử cho mọi nhánh}
\NormalTok{→ cài đặt bằng Python}
\NormalTok{→ chạy ca thường, ca biên và ca không hợp lệ}
\NormalTok{→ mô phỏng khi kết quả bất thường}
\end{Highlighting}
\end{Shaded}

Một chu trình thực hành hữu ích:

\begin{quote}
\textbf{Dự đoán → Mô phỏng → Chạy → So sánh → Giải thích → Sửa đổi →
Kiểm tra biên → Luyện tập}
\end{quote}

Thói quen quan trọng nhất của chương này là coi mỗi quyết định là một
cấu trúc cần được \textbf{thiết kế và kiểm thử}, không chỉ đơn thuần là
gõ vài câu lệnh \texttt{if}.

\end{document}
